\documentclass[14pt,a4paper]{extarticle}

\usepackage{fontspec}
\usepackage{polyglossia}
\setdefaultlanguage{russian}
\setotherlanguage{english}
\setmainfont{Times New Roman}
\newfontfamily\cyrillicfont{Times New Roman}[Script=Cyrillic]
\newfontfamily\cyrillicfonttt{Courier New}[Script=Cyrillic]

\usepackage{geometry}
\geometry{left=30mm,right=10mm,top=20mm,bottom=20mm}
\usepackage{setspace}
\onehalfspacing
\usepackage{indentfirst}
\setlength{\parindent}{1.25cm}
\usepackage{array}
\usepackage{longtable}
\usepackage{booktabs}
\usepackage{tabularx}
\usepackage{float}
\usepackage{amsmath}
\usepackage{amssymb}
\usepackage{graphicx}
\usepackage{enumitem}
\usepackage{xcolor}
\usepackage[hidelinks]{hyperref}
\usepackage{titlesec}
\usepackage{tocloft}

\titleformat{\section}{\normalfont\bfseries\centering\MakeUppercase}{\thesection}{1em}{}
\titleformat{name=\section,numberless}{\normalfont\bfseries\centering\MakeUppercase}{}{0pt}{}
\titleformat{\subsection}{\normalfont\bfseries}{\thesubsection}{1em}{}
\titleformat{\subsubsection}{\normalfont\bfseries}{\thesubsubsection}{1em}{}
\titlespacing*{\section}{0pt}{0pt}{12pt}
\titlespacing*{\subsection}{0pt}{12pt}{6pt}
\titlespacing*{\subsubsection}{0pt}{10pt}{4pt}

\renewcommand{\cftsecleader}{\cftdotfill{\cftdotsep}}
\setlist{nosep,leftmargin=1.25cm}
\setlength{\emergencystretch}{3em}
\sloppy

\newcommand{\todoBox}[1]{%
  \begin{center}
  \fbox{\begin{minipage}{0.94\textwidth}
  \textbf{Нужно добавить перед сдачей.} #1
  \end{minipage}}
  \end{center}
}

\newcommand{\source}[1]{[#1]}

\begin{document}

\begin{titlepage}
\begin{center}
ПРАВИТЕЛЬСТВО РОССИЙСКОЙ ФЕДЕРАЦИИ\\
НАЦИОНАЛЬНЫЙ ИССЛЕДОВАТЕЛЬСКИЙ УНИВЕРСИТЕТ\\
«ВЫСШАЯ ШКОЛА ЭКОНОМИКИ»\\
Факультет компьютерных наук\\
Образовательная программа «Программная инженерия»\\
\vspace{0.8cm}
УДК 004.023\\
\vspace{2.2cm}
\textbf{ОТЧЕТ}\\
по исследовательскому курсовому проекту\\
\vspace{0.8cm}
на тему\\
\textbf{«Оптимизация алгоритмов маршрутизации на примере задачи нескольких коммивояжеров»}\\
\vspace{0.8cm}
по направлению подготовки бакалавров 09.03.04 «Программная инженерия»
\end{center}

\vfill

\begin{flushright}
Выполнил\\
студент группы БПИ246\\
Д. Д. Купцов\\
\vspace{0.7cm}
Научный руководитель\\
Ермаков Андрей Максимович
\end{flushright}

\vfill

\begin{center}
Москва 2026
\end{center}
\end{titlepage}

\setcounter{page}{2}

\section*{РЕФЕРАТ}

Отчет содержит описание исследовательского курсового проекта, посвященного задаче множественного коммивояжера и практическому сравнению эвристических алгоритмов маршрутизации. В работе рассматривается постановка mTSP с одним депо и целевой функцией MINSUM, при которой требуется минимизировать суммарную длину всех маршрутов. Объектом исследования являются процессы маршрутизации и распределения точек между несколькими агентами, предметом исследования являются эвристические и метаэвристические методы решения mTSP при ограниченном времени вычислений.

Цель работы --- разработать и экспериментально оценить масштабируемые эвристические методы оптимизации маршрутов для задачи множественного коммивояжера. Для достижения цели были изучены классические подходы к TSP и mTSP, реализованы базовые алгоритмы и семейство LKH-подобных оберток, подготовлен экспериментальный контур на C++ и Python, проведены сравнения с OR-Tools, преобразованием TSP $\rightarrow$ mTSP и внешним решателем LKH-3. Полученные результаты показывают, что LKH-подобные методы существенно превосходят простые конструктивные эвристики по качеству маршрутов: на равномерных инстансах средний относительный gap к OR-Tools снижается с~114.7\% у~\texttt{rand+nn} до~14.5\% у~\texttt{lkh-wrapper-v2}, а попарное сравнение \texttt{v2} vs \texttt{grasp} на~120 инстансах даёт 93~победы. Однако при сравнении с полноформатным LKH-3 возникает компромисс: собственная версия~\texttt{v12} в~6.2~раза быстрее на инстансах $n=5000$, но уступает по~MINSUM около~21\%; при этом~\texttt{v12} строит маршруты с заметно меньшим разбалансом (1.0--2.7~против~1.7--5.0 у~LKH-3), что согласуется с теоретическими работами по связи MINSUM- и MINMAX-целей в режиме большого числа агентов. Дальнейшее развитие проекта (\texttt{v13}--\texttt{v21}) переориентировано на архитектуру Adaptive Large Neighborhood Search с интегрированным Simulated Annealing и опциональным Parallel Tempering; на этой архитектуре проведён formal variance audit на flagship-инстансах ($n=10\,000$, $50\,000$, $100\,000$, по~10 seeds на каждый), показавший относительное стандартное отклонение MINSUM в пределах 0.17--0.59\% и идентифицировавший super-linear scaling-overhead $\approx 2.14\times$ при росте $n$ как ключевое архитектурное ограничение. Отдельным экспериментом для сценариев с большим числом агентов ($m=100$) был добавлен FILO2-вдохновлённый solver \texttt{lkh\_v21\_minsum\_cap} с capacity-aware repair: на четырёх high-m инстансах он даёт парный MINSUM-выигрыш в~3.24\,\% в среднем (16~из~20~seeds улучшены, Wilcoxon $p=0.002$), сокращая разрыв до FILO2 на flagship-инстансе с~$\sim$30\,\% до~$\sim$15.7\,\%.

\textbf{Ключевые слова:} задача коммивояжера, множественный коммивояжер, mTSP, маршрутизация, эвристические алгоритмы, Lin--Kernighan--Helsgaun, локальный поиск, MINSUM, вычислительный эксперимент.

% TODO: после финальной компиляции PDF уточнить число страниц, таблиц, рисунков, источников и приложений в строке реферата выше.

\newpage
\tableofcontents
\newpage

\section*{ОПРЕДЕЛЕНИЯ, ОБОЗНАЧЕНИЯ И СОКРАЩЕНИЯ}
\addcontentsline{toc}{section}{ОПРЕДЕЛЕНИЯ, ОБОЗНАЧЕНИЯ И СОКРАЩЕНИЯ}

\begin{longtable}{>{\raggedright\arraybackslash}p{0.28\textwidth}>{\raggedright\arraybackslash}p{0.66\textwidth}}
\toprule
\textbf{Термин} & \textbf{Определение}\\
\midrule
\endhead
TSP & Travelling Salesman Problem, задача коммивояжера: требуется построить кратчайший замкнутый маршрут, проходящий через все заданные вершины ровно один раз.\\
mTSP & Multiple Travelling Salesman Problem, задача множественного коммивояжера: требуется построить несколько маршрутов, начинающихся и заканчивающихся в общем депо, так чтобы каждая клиентская вершина была посещена ровно одним маршрутом.\\
Депо & Специальная вершина, из которой начинаются и в которой завершаются маршруты всех агентов.\\
Агент, коммивояжер & Исполнитель, которому назначается отдельный маршрут в решении mTSP.\\
MINSUM & Целевая функция mTSP, равная суммарной длине всех маршрутов. В данной работе меньшая MINSUM означает более качественное решение.\\
MINMAX & Альтернативная целевая функция, в которой минимизируется длина самого длинного маршрута. В работе она используется как вспомогательный контекст, но не как основной критерий.\\
2-opt & Локальное улучшение маршрута, при котором удаляются два ребра и фрагмент маршрута разворачивается, если это уменьшает длину.\\
k-opt & Обобщение 2-opt: из маршрута удаляется $k$ ребер, после чего фрагменты соединяются новым способом.\\
LKH & Lin--Kernighan--Helsgaun, семейство высокоэффективных эвристик для TSP, основанных на переменной глубине локального поиска и кандидатных ребрах.\\
Candidate set & Ограниченный набор перспективных соседей вершины, используемый для ускорения локального поиска и отказа от полного перебора всех пар вершин.\\
GRASP & Greedy Randomized Adaptive Search Procedure, жадно-рандомизированная процедура построения и улучшения решений.\\
OR-Tools & Библиотека Google для решения задач оптимизации, в работе используется как внешний эвристический ориентир для маршрутизации.\\
FILO2 & Эвристический подход для крупномасштабных задач маршрутизации, используемый как внешний ориентир по идеям масштабирования.\\
Benchmark & Набор тестовых экземпляров задачи, на котором сравниваются алгоритмы.\\
Relative gap & Относительное отклонение решения от выбранного ориентира: $100 \cdot (A - B) / B$, где $A$ --- значение исследуемого алгоритма, $B$ --- значение базового решения.\\
\bottomrule
\end{longtable}

\newpage
\section*{ВВЕДЕНИЕ}
\addcontentsline{toc}{section}{ВВЕДЕНИЕ}

Задачи маршрутизации относятся к числу наиболее практически значимых задач комбинаторной оптимизации. Они возникают в логистике, транспортном планировании, управлении цепями поставок, распределении сервисных бригад, проектировании производственных процессов и других областях, где требуется упорядочить перемещение между множеством точек. Классическим математическим ядром таких задач является задача коммивояжера, в которой необходимо найти маршрут минимальной длины, проходящий через все вершины ровно один раз и возвращающийся в исходную вершину \source{1,2}.

Несмотря на простую формулировку, TSP относится к NP-трудным задачам. Это означает, что при росте числа вершин полный перебор решений становится неприемлемым уже на сравнительно малых размерах. В прикладных постановках ситуация обычно еще сложнее: маршрут строится не для одного исполнителя, а для нескольких агентов, которые стартуют из общего депо и должны совместно обслужить набор клиентов. Такая постановка называется задачей множественного коммивояжера, или mTSP. В отличие от TSP, здесь требуется не только определить порядок обхода вершин, но и распределить вершины между маршрутами \source{4}.

В данной работе рассматривается вариант mTSP с одним депо и целевой функцией MINSUM. Для практических задач такая постановка естественна: если несколько машин или исполнителей обслуживают один и тот же набор точек, то суммарная длина маршрутов может быть связана с расходом топлива, временем работы, стоимостью обслуживания или суммарной нагрузкой на систему. При этом получение точного оптимума для задач средней и большой размерности обычно не является реалистичной целью. Более важным становится поиск достаточно качественного приближенного решения в заранее заданный вычислительный бюджет.

Актуальность исследования определяется двумя обстоятельствами. Во-первых, mTSP является базовой моделью для более сложных задач маршрутизации, включая варианты Vehicle Routing Problem; современный обзор~\source{35} насчитывает несколько десятков прикладных направлений, в которых mTSP используется как ядро. Во-вторых, современные эвристические методы показывают, что сильное качество решения часто достигается не за счет полного перебора, а за счет аккуратного ограничения области поиска: кандидатных ребер, локальных перестановок, кластеризации, разреженных соседств и процедур повторного улучшения \source{3,7,8,10,28,29}. Поэтому практический интерес представляет не только выбор известного решателя, но и анализ того, какие инженерные решения делают алгоритм масштабируемым.

Объектом исследования являются процессы маршрутизации и распределения маршрутов между несколькими агентами в задачах комбинаторной оптимизации. Предметом исследования являются эвристические и метаэвристические алгоритмы решения mTSP, а также их влияние на качество маршрутов и время вычислений при росте размерности задачи.

Цель работы --- разработать и проанализировать эффективные эвристические методы оптимизации маршрутов для задачи множественного коммивояжера. Для достижения цели в работе решаются следующие задачи:

\begin{enumerate}
  \item изучить постановку TSP и mTSP, основные целевые функции и ограничения;
  \item проанализировать существующие эвристические подходы, включая локальный поиск, GRASP, LKH и методы построения candidate sets;
  \item реализовать единый экспериментальный контур для запуска алгоритмов и проверки корректности решений;
  \item подготовить несколько семейств синтетических mTSP-инстансов, отражающих разные геометрические сценарии;
  \item сравнить базовые эвристики, собственные LKH-подобные версии и внешние ориентиры по качеству, времени и валидности решений;
  \item сформулировать выводы о применимости разработанных подходов и ограничениях текущей реализации.
\end{enumerate}

В качестве рабочей гипотезы принимается следующее положение: эвристики, использующие ограниченные множества кандидатных ребер, локальный поиск и LKH-подобные процедуры улучшения, позволяют получать более качественные решения mTSP по MINSUM, чем простые конструктивные алгоритмы, при сопоставимых или приемлемых вычислительных затратах. Дополнительно проверяется более осторожная инженерная гипотеза: быстрые собственные версии могут быть полезны при жестком временном бюджете, даже если по качеству они уступают полноформатному внешнему LKH-3.

Методологическая основа работы включает анализ научной литературы, математическую формализацию задачи, программную реализацию алгоритмов, вычислительный эксперимент, валидацию маршрутов и сравнительный анализ результатов. Основные эксперименты проводятся на синтетических инстансах с разными распределениями точек: равномерным, кластерным, кластерным со смещенным депо, смешанным с выбросами и стрессовыми сценариями с большим числом агентов.

Все эксперименты проводились на одной машине. Аппаратная конфигурация: процессор Intel Core Ultra~7 165U (12 физических ядер, 14 логических потоков), оперативная память 32~ГБ, операционная система Windows~11~Pro (версия 10.0.26200). Сборка выполнена компилятором Clang~20.1.2 через CMake~4.2.1 в режиме Release с флагом оптимизации \texttt{-O2} и включенным OpenMP для параллельного polish-прохода. Число потоков OpenMP задается параметром \texttt{--threads} (по умолчанию соответствует числу агентов~$m$).

\newpage
\section{АНАЛИЗ ПРЕДМЕТНОЙ ОБЛАСТИ}

\subsection{Классическая задача коммивояжера}

Задача коммивояжера формулируется следующим образом. Пусть задано множество вершин $V=\{0,1,\ldots,n-1\}$ и функция расстояния $d(i,j)$ между каждой парой вершин. Требуется найти перестановку вершин, задающую замкнутый маршрут минимальной длины:
\begin{equation}
  \min_{\pi} \left(d(\pi_n,\pi_1) + \sum_{i=1}^{n-1} d(\pi_i,\pi_{i+1})\right).
\end{equation}
В евклидовой постановке вершины задаются координатами на плоскости, а расстояние вычисляется как длина отрезка между точками. Такая постановка удобна для экспериментов, поскольку позволяет генерировать тестовые данные контролируемым образом и сохраняет связь с прикладными задачами маршрутизации.

Точные методы решения TSP имеют большую теоретическую ценность, но для крупных экземпляров используются ограниченно. В ранних работах были предложены процедуры ветвей и границ, динамического программирования и другие точные подходы, однако рост пространства решений делает их дорогостоящими \source{1}. Поэтому в практических задачах широко применяются эвристики. Уже классический алгоритм Lin--Kernighan показал, что локальный поиск переменной глубины способен давать оптимальные или близкие к оптимальным решения на широком классе тестов \source{1}. Позднее Helsgaun развил этот подход в LKH, усилив его использованием кандидатных ребер и более сложной стратегии поиска \source{2}.

\subsection{Задача множественного коммивояжера}

mTSP является естественным обобщением TSP. Пусть задано $m$ агентов и общее депо $0$. Решение представляет собой набор маршрутов:
\begin{equation}
  R_s = (0, v^s_1, v^s_2,\ldots,v^s_{k_s}, 0), \quad s=1,\ldots,m,
\end{equation}
где каждая вершина из $V \setminus \{0\}$ встречается ровно в одном маршруте. В варианте MINSUM целевая функция имеет вид:
\begin{equation}
  F_{\mathrm{sum}}(R_1,\ldots,R_m)=
  \sum_{s=1}^{m}\sum_{i=0}^{k_s} d(v^s_i, v^s_{i+1}),
  \quad v^s_0=v^s_{k_s+1}=0.
\end{equation}

В литературе также часто рассматривается MINMAX, где минимизируется длина самого длинного маршрута. Эти цели отражают разные практические смыслы. MINSUM ориентирован на суммарную стоимость обслуживания, а MINMAX --- на балансировку нагрузки между агентами. В данной работе основной критерий --- MINSUM, но баланс маршрутов дополнительно анализируется через отношение длины самого длинного маршрута к средней длине маршрута. Связь между MINSUM и MINMAX-целями формально исследована в работе \source{20}: на наихудших инстансах разрыв между MINSUM- и MINMAX-оптимумами линейно растёт по числу агентов~$m$, однако в типичных режимах (равномерное распределение, большое~$m$) балансированные маршруты дают суммарную длину, близкую к MINSUM-оптимуму~\source{21,22}. Эмпирический обзор \source{23} показывает, что сокращение разброса длин маршрутов на~40\% достигается ценой увеличения суммарной длины лишь на~2--3\%. Это служит мотивацией к интерпретации наблюдаемого баланса \texttt{lkh-wrapper-v12} в разделе~5.6.

Обзор Bektaş~\source{4} показывает, что mTSP имеет несколько эквивалентных и близких формулировок, а также может решаться через преобразования к TSP, целочисленное программирование и эвристические методы. Современное обновление этого обзора~\source{35} предлагает таксономию вариантов mTSP по числу депо, целевой функции, ограничениям и методам решения и охватывает работы 2006--2021~годов. Для задач большой размерности особенно важны эвристики, поскольку даже если точный метод может решить малый пример, он не дает приемлемого времени на крупных данных.

\subsection{Эвристические и метаэвристические подходы}

Простейший класс методов составляют конструктивные эвристики. Они строят решение последовательно, выбирая следующую вершину по локальному правилу, например по ближайшему соседу. Такие методы очень быстры, но часто принимают ранние решения, которые затем трудно исправить. В проекте к этой группе относится алгоритм \texttt{rand+nn}: он распределяет вершины между маршрутами с элементом случайности и ближайшего соседа.

Второй класс методов основан на локальном поиске. Для каждого маршрута можно применять 2-opt: если замена двух ребер уменьшает длину маршрута, соответствующий фрагмент разворачивается. В проекте этот подход представлен алгоритмом \texttt{2opt+greed}. Он сильнее чистого ближайшего соседа, но все еще ограничен, потому что улучшает уже построенные маршруты и не всегда хорошо перераспределяет вершины между агентами.

Третий класс составляют рандомизированные многостартовые процедуры. GRASP строит несколько решений с использованием ограниченного списка лучших локальных кандидатов, затем улучшает каждое решение и выбирает лучшее \source{9}. В реализации проекта \texttt{grasp} использует randomized restricted candidate list, внутримаршрутный 2-opt и межмаршрутные обмены. Такой подход заметно повышает качество, но его время растет с числом итераций и размером инстанса.

Отдельное место занимают LKH-подобные методы. Их сила состоит в том, что поиск не рассматривает все возможные ребра, а концентрируется на малом наборе перспективных соседей. Для TSP эта идея поддержана классическими работами Lin--Kernighan и Helsgaun \source{1,2}. Для mTSP LKH-3 использует преобразование задач маршрутизации к стандартному симметричному TSP и штрафные функции для ограничений \source{3}. В частности, для mTSP с одним депо возможно добавить копии депо и получить структуру, распадающуюся на несколько маршрутов.

В современной литературе по VRP доминирующее положение занимают два класса метаэвристик, которые не реализованы в данном проекте напрямую, но задают методологический контекст. Первый класс --- Adaptive Large Neighborhood Search~(ALNS) Ropke и Pisinger~\source{24,25}, основанный на чередовании destroy- и repair-операторов с адаптивным выбором их веса. Второй класс --- Hybrid Genetic Search~(HGS) Vidal и соавт.~\source{26,27}, объединяющий популяционный поиск, локальное улучшение и Split-алгоритм расстановки разделителей маршрутов. Развитие ALNS в работе Christiaens и Vanden~Berghe~\source{28} (Slack Induction by String Removals,~SISR) показало, что простой и единственный destroy-оператор удаления непрерывных строк клиентов из пространственно близких маршрутов обеспечивает SOTA-качество на CVRP/VRPTW/MDVRP. Эти методы важны как ориентир для дальнейшего развития~\texttt{lkh-wrapper}: они показывают, что современный VRP-солвер требует не только candidate-based локального поиска, но и более богатой destroy/repair- или популяционной инфраструктуры.

\subsection{Масштабирование и candidate sets}

Ключевая инженерная проблема крупных TSP/mTSP-инстансов --- невозможность рассматривать полный граф. Если для каждой вершины проверять все остальные вершины, построение соседств требует порядка $O(n^2)$ сравнений. Для $n=100000$ это уже непрактично. Поэтому современные эвристики строят разреженный граф кандидатов.

В работах Taillard и Helsgaun по POPMUSIC показано, что качественные candidate sets можно строить с эмпирической сложностью ниже квадратичной даже для очень крупных TSP-инстансов \source{7}. Похожую идею используют крупномасштабные VRP-решатели: они ограничивают локальный поиск гранулярными соседствами, динамическими областями изменений и кэшами затронутых вершин \source{8}. Концепция гранулярных окрестностей восходит к работе Toth и~Vigo~\source{29}, где предложено исключать из рассмотрения рёбра длины более~$\beta \cdot \bar{c}$, где $\bar{c}$~--- средняя длина ребра в текущем решении, а~$\beta$~--- настраиваемый коэффициент. Это сужает окрестность на порядок при минимальной потере качества и лежит в основе candidate-генерации в LKH, neighbor-lists в OR-Tools и большинстве современных эвристик для VRP~\source{30}. Для данной работы это важно не как прямое заимствование готового алгоритма, а как подтверждение общего принципа: масштабируемость достигается через сужение пространства ходов и отказ от полного перебора.

В собственных версиях \texttt{lkh-wrapper} этот принцип реализован через геометрические candidate sets, пространственную сетку, кэширование расстояний, ограниченный 2-opt и time-budget. В более поздних версиях добавлены кластерные начальные маршруты, процедуры перераспределения блоков между маршрутами и отдельные бюджеты на построение первого решения и последующее улучшение.

\newpage
\section{ПОДРОБНЫЙ РАЗБОР ЭВРИСТИК MTSP}

В данной главе отдельно рассматриваются эвристики, которые используются в проекте или служат для него методологическими ориентирами. Такое разделение важно по двум причинам. Во-первых, итоговое качество решения mTSP определяется не только формальной постановкой, но и тем, как именно построено начальное распределение вершин, какие локальные ходы разрешены и насколько сильно ограничено пространство поиска. Во-вторых, разные алгоритмы полезны в разных режимах: один метод удобен как быстрый baseline, другой дает лучшее качество на малых инстансах, третий показывает, какие инженерные идеи нужны для масштабирования до десятков и сотен тысяч вершин.

Основная целевая функция в проекте --- MINSUM. Поэтому при сравнении эвристик главный вопрос формулируется так: насколько эффективно алгоритм уменьшает суммарную длину всех маршрутов при сохранении валидности решения. При этом отдельно учитывается баланс маршрутов, поскольку слишком длинный отдельный маршрут может быть нежелателен даже тогда, когда суммарная длина выглядит приемлемой.

\begin{figure}[H]
\centering
\fbox{\begin{minipage}{0.92\textwidth}
\centering
конструктивные эвристики (\texttt{rand+nn, 2opt+greed}) $\;\rightarrow\;$ многостартовый локальный поиск (\texttt{grasp}) $\;\rightarrow\;$ LKH-подобные версии (\texttt{lkh-wrapper-v1\ldots v12}) $\;\rightarrow\;$ внешние ориентиры (LKH-3, FILO2, POPMUSIC)
\end{minipage}}
\caption{Иерархия эвристик: от простых конструктивных методов к крупномасштабным ориентирам из литературы}
\end{figure}

\subsection{\texttt{mtsp\_rand\_nn\_solver}: случайное разбиение и nearest neighbor}

Решатель \texttt{mtsp\_rand\_nn\_solver} зарегистрирован в проекте под именем \texttt{rand+nn}. Это самый простой внутренний baseline, который нужен не столько для получения сильного результата, сколько для проверки экспериментального контура и оценки нижнего уровня качества простых конструктивных решений.

Алгоритм работает в два этапа. Сначала все клиентские вершины, то есть вершины $1,\ldots,n-1$, перемешиваются с использованием фиксируемого seed. Затем перемешанный список распределяется между коммивояжерами по кругу. После этого для каждой полученной группы строится маршрут ближайшего соседа: текущая вершина последовательно соединяется с ближайшей еще не посещенной вершиной из той же группы.

Сильная сторона этого подхода --- очень малое время работы и воспроизводимость при заданном seed. Он удобен как технический ориентир: если более сложная эвристика проигрывает \texttt{rand+nn}, значит в ней, скорее всего, есть ошибка реализации или неудачная настройка параметров. Кроме того, случайное начальное разбиение позволяет быстро получать несколько разных решений без изменения основной логики.

Главный недостаток метода состоит в том, что распределение клиентов между маршрутами никак не связано с геометрией задачи. Две близкие точки могут оказаться у разных агентов, а один маршрут может получить группу, которая геометрически плохо собирается в короткий цикл. Локальный nearest neighbor внутри группы не исправляет ошибку назначения, потому что не переносит вершины между маршрутами. Поэтому \texttt{rand+nn} ожидаемо дает худшее качество среди реализованных внутренних методов, но остается полезным как базовая точка сравнения.

\subsection{\texttt{mtsp\_greedy\_2opt\_solver}: жадное построение и 2-opt}

Решатель \texttt{mtsp\_greedy\_2opt\_solver} зарегистрирован под именем \texttt{2opt+greed}. Он усиливает простую конструктивную схему за счет более осмысленного выбора следующей вершины и последующего локального улучшения каждого маршрута.

Построение начинается с $m$ пустых маршрутов, каждый из которых стартует в депо. Далее алгоритм циклически проходит по коммивояжерам. Для текущего коммивояжера выбирается ближайшая еще не посещенная вершина относительно конца его текущего маршрута. Выбранная вершина добавляется в маршрут, помечается как посещенная, и процесс продолжается, пока все клиентские вершины не будут распределены. После закрытия маршрутов возвратом в депо для каждого маршрута отдельно запускается процедура 2-opt.

2-opt улучшает порядок обхода внутри одного маршрута. Если две дуги маршрута можно заменить двумя другими дугами так, что цикл станет короче, соответствующий фрагмент маршрута разворачивается. В контексте данного решателя это означает, что жадное распределение вершин остается прежним, но порядок посещения внутри каждого маршрута становится более аккуратным. Поэтому \texttt{2opt+greed} обычно заметно лучше \texttt{rand+nn}: он устраняет часть пересечений и явно неудачных участков внутри маршрутов.

Ограничение метода связано с тем, что он почти не пересматривает ранние решения. Если вершина была назначена одному агенту, 2-opt не передаст ее другому агенту, даже если такой перенос уменьшил бы MINSUM. Кроме того, жадный выбор ближайшей вершины может создавать локально хорошие, но глобально неудачные продолжения. По этой причине \texttt{2opt+greed} является сильнее простого baseline, но все еще уступает многостартовым и LKH-подобным методам.

\begin{figure}[H]
\centering
\fbox{\begin{minipage}{0.92\textwidth}
\centering
Жадное назначение ближайшей вершины каждому агенту $\rightarrow$ закрытие маршрутов в депо $\rightarrow$ внутримаршрутный 2-opt $\rightarrow$ проверка валидности и расчет MINSUM
\end{minipage}}
\caption{Логика эвристики \texttt{2opt+greed}}
\end{figure}

% TODO (опционально): заменить текстовую схему выше на PNG-визуализацию маршрутов до/после 2-opt на малом инстансе.

\subsection{\texttt{mtsp\_grasp\_solver}: GRASP, RCL и межмаршрутные обмены}

\texttt{mtsp\_grasp\_solver} реализует GRASP-подход и зарегистрирован как \texttt{grasp}. В отличие от \texttt{2opt+greed}, он не строит единственное детерминированное решение. Алгоритм многократно создает кандидатные решения с элементом случайности, улучшает их и сохраняет лучшее найденное значение MINSUM.

На этапе построения для каждого агента формируется список доступных городов, отсортированный по расстоянию от текущей вершины маршрута. Вместо того чтобы всегда брать ближайший город, алгоритм выбирает случайную вершину из restricted candidate list --- списка из нескольких лучших локальных кандидатов. Размер этого списка задается параметром \texttt{rcl}. Если \texttt{rcl}=1, поведение становится близким к жадному выбору; при большем \texttt{rcl} поиск получает больше разнообразия, но может чаще делать слабые локальные шаги.

После построения каждого маршрута применяется 2-opt. Затем запускается межмаршрутная фаза: алгоритм перебирает пары маршрутов и пробует поменять местами внутренние вершины. Если обмен уменьшает MINSUM, он принимается, а затронутые маршруты дополнительно улучшаются 2-opt. Этот шаг принципиально отличает \texttt{grasp} от чистого \texttt{2opt+greed}: теперь алгоритм может частично исправлять неудачное распределение вершин между агентами.

Преимущество GRASP состоит в сочетании трех механизмов: рандомизации, локального улучшения и выбора лучшего решения из нескольких запусков. Такой подход устойчивее к локальным ошибкам начального построения и обычно дает более качественные маршруты, чем чистые конструктивные методы. Недостаток --- рост времени работы. На каждом старте требуется сортировать локальных кандидатов, выполнять 2-opt и проверять межмаршрутные обмены. При увеличении числа итераций качество может улучшаться, но вычислительная стоимость растет почти прямо вместе с числом попыток.

\begin{longtable}{>{\raggedright\arraybackslash}p{0.26\textwidth}>{\raggedright\arraybackslash}p{0.28\textwidth}>{\raggedright\arraybackslash}p{0.36\textwidth}}
\toprule
\textbf{Параметр} & \textbf{Смысл} & \textbf{Влияние на результат}\\
\midrule
\endhead
\texttt{iters} & число независимых построений решения & повышает шанс найти более короткий MINSUM, но линейно увеличивает время работы\\
\texttt{rcl} & размер restricted candidate list & управляет балансом между жадностью и разнообразием поиска\\
\texttt{seed} & начальное состояние генератора случайных чисел & делает стохастический запуск воспроизводимым\\
\bottomrule
\end{longtable}

\begin{figure}[H]
\centering
\fbox{\begin{minipage}{0.92\textwidth}
\centering
все незанятые вершины, отсортированные по $d(\text{текущая}, v_i)$: $[v_{(1)}, v_{(2)}, \ldots, v_{(n)}]$\\[4pt]
$\underbrace{v_{(1)}, v_{(2)}, \ldots, v_{(\texttt{rcl})}}_{\text{restricted candidate list}}$ $\leftarrow$ случайный выбор из этого подсписка
\end{minipage}}
\caption{Принцип restricted candidate list в GRASP: не всегда брать ближайшую вершину, а выбирать случайную из топ-$k$ кандидатов}
\end{figure}

\subsection{LKH и LKH-3 как внешний ориентир качества}

Lin--Kernighan и LKH относятся к более сильному классу локального поиска переменной глубины. В обычном 2-opt фиксируется простой тип хода: заменить две дуги. В Lin--Kernighan логика гибче: алгоритм строит последовательность замен ребер и продолжает ее, пока есть надежда получить выигрыш. Helsgaun в LKH усилил этот подход candidate sets, специальными правилами выбора ребер и настройками, которые позволяют не просматривать полный граф \source{1,2}.

Для mTSP особенно важен LKH-3. Он расширяет LKH на ряд ограниченных задач маршрутизации, включая варианты TSP, VRP и mTSP. Общая идея состоит в сведении исходной задачи к стандартному симметричному TSP с дополнительными структурами и штрафными функциями \source{3}. Для задачи нескольких коммивояжеров это позволяет описывать число агентов, депо и целевую функцию, например MINSUM. В практическом смысле LKH-3 выступает сильным внешним baseline: если собственная эвристика приближается к нему по MINSUM, это серьезный признак качества.

В проекте LKH-3 используется не как основная внутренняя реализация, а как внешний ориентир. Такой выбор оправдан: полноценный LKH-3 уже содержит большое число инженерных и алгоритмических оптимизаций, поэтому честнее сравнивать собственные версии с ним как с сильным решателем, а не пытаться представить его как часть собственной разработки. Сравнение с LKH-3 также помогает понять, где именно внутренним \texttt{lkh-wrapper} еще не хватает глубины поиска.

Ограничение LKH-3 в рамках проекта связано с интеграцией и управлением экспериментами. Внешний запуск требует подготовки файлов задачи и параметров, запуска отдельного исполняемого файла и разбора результата. Кроме того, сильный внешний решатель не всегда удобно адаптировать под исследовательские изменения, например под собственные ограничения, альтернативные метрики баланса или специальные бюджеты времени. Поэтому в работе параллельно развиваются собственные LKH-подобные версии.

\begin{figure}[H]
\centering
\fbox{\begin{minipage}{0.92\textwidth}
\centering
mTSP с $m$ агентами и депо $0$\\[4pt]
$\downarrow$ добавить $m{-}1$ копий депо, задать penalty function для числа маршрутов\\[4pt]
стандартный симметричный TSP с ограничениями\\[4pt]
$\downarrow$ LKH-поиск переменной глубины + candidate sets\\[4pt]
решение TSP $\;\rightarrow\;$ декодирование в $m$ маршрутов mTSP
\end{minipage}}
\caption{Общая схема преобразования mTSP к TSP в LKH-3 \source{3}}
\end{figure}

\subsection{Собственные LKH-подобные версии \texttt{lkh-wrapper}}

Семейство \texttt{lkh-wrapper} в проекте не является прямой копией LKH-3. Это набор собственных эвристических версий, которые используют идеи LKH-подхода в упрощенном и более контролируемом виде. Главная инженерная цель этих версий --- получить валидное mTSP-решение за ограниченное время и улучшать его без полного перебора всех возможных ребер.

Ранние версии \texttt{v1}--\texttt{v3} использовались как переход от простых baseline-алгоритмов к candidate-based логике. Они показывают, что даже умеренное ограничение поиска перспективными соседями может дать большой выигрыш по MINSUM. Однако эти версии еще рассчитаны скорее на малые и средние инстансы.

Версии \texttt{v4}--\texttt{v7} были направлены на масштабирование. В них появились геометрические candidate sets, пространственная сетка, кэширование расстояний и регулярная проверка time-budget. Это важный шаг: на больших инстансах нельзя хранить и просматривать полную матрицу расстояний как основной рабочий объект. Алгоритм должен быстро получать только те расстояния и соседства, которые реально участвуют в поиске.

Поздняя версия \texttt{lkh-wrapper-v12} ориентирована на MINSUM и быстрый возврат первого валидного решения. В ней используются отдельные бюджеты на построение начального решения и на последующее улучшение, а также идеи cluster-aware построения, savings-подхода и межмаршрутного перераспределения. По смыслу эта версия ближе всего к практическому решателю: она не просто улучшает один маршрут, а пытается управлять распределением вершин между агентами.

Сильная сторона собственных LKH-подобных версий --- прозрачность и управляемость. Их можно запускать внутри общего экспериментального контура, менять параметры, измерять влияние отдельных эвристических блоков и адаптировать под конкретную постановку. Слабая сторона --- меньшая глубина поиска по сравнению с полноценным LKH-3. Результаты экспериментов показывают, что \texttt{v12} может быть быстрее внешнего LKH-3, но обычно уступает ему по MINSUM. Это естественный компромисс между скоростью, простотой интеграции и качеством решения.

\begin{figure}[H]
\centering
\fbox{\begin{minipage}{0.92\textwidth}
\centering
candidate sets $\rightarrow$ начальные маршруты $\rightarrow$ внутримаршрутное улучшение $\rightarrow$ межмаршрутные переносы/обмены $\rightarrow$ контроль time-budget $\rightarrow$ лучшее валидное решение
\end{minipage}}
\caption{Обобщенная схема LKH-подобных версий проекта}
\end{figure}

\subsection{FILO2 и крупномасштабные VRP-эвристики}

FILO2 рассматривается в работе как важный внешний ориентир по масштабированию, хотя сам он относится к крупномасштабным задачам маршрутизации транспорта, а не к чистой базовой постановке mTSP. Его ценность для данного проекта состоит в том, что статья Accorsi и Vigo показывает практический путь к решению очень больших маршрутизационных инстансов: алгоритм не пытается проверять все возможные перемещения, а концентрирует локальный поиск вокруг небольшой области, где изменение действительно может дать эффект \source{8}.

В FILO2 используются несколько идей, полезных для понимания дальнейшего развития проекта. Granular neighborhoods ограничивают число рассматриваемых соседей и тем самым уменьшают стоимость локального поиска. Static move descriptors позволяют заранее описывать типы перемещений и быстрее применять их к решению. Sequential search организует улучшения последовательно, без тяжелого полного перебора всех комбинаций. Selective vertex caching хранит информацию о вершинах, которые могли быть затронуты последними изменениями, и не тратит время на области, где улучшение маловероятно.

Для mTSP эти идеи можно интерпретировать следующим образом. Если вершина была перенесена между маршрутами или внутри маршрута изменился фрагмент, нет необходимости заново анализировать весь набор клиентов. Достаточно проверить локальную окрестность измененных вершин, соседние участки маршрутов и несколько кандидатов на межмаршрутный перенос. Такой принцип хорошо согласуется с уже реализованными candidate sets и time-budget в \texttt{lkh-wrapper}.

При этом FILO2 нельзя без оговорок считать готовым mTSP-решателем для данной работы. В исходной статье речь идет о более общей VRP-среде, где присутствуют ограничения и цели, отличающиеся от простой MINSUM-постановки mTSP. Поэтому в отчете FILO2 используется как источник архитектурных идей: как ограничивать область поиска, как кэшировать затронутые вершины и как строить эвристику, которая остается работоспособной на очень больших данных.

\begin{figure}[H]
\centering
\fbox{\begin{minipage}{0.92\textwidth}
\centering
после локального изменения (перенос или замена вершин):\\[4pt]
все вершины $\rightarrow$ \textit{selective vertex cache}: только вершины вблизи изменения\\[4pt]
$\downarrow$ \textit{granular neighborhood}: ограниченный список кандидатов для каждой вершины\\[4pt]
$\downarrow$ локальный поиск только в пострадавшей области\\[4pt]
$\Rightarrow$ линейная по числу затронутых вершин стоимость одного улучшающего шага
\end{minipage}}
\caption{Идея selective vertex caching и granular neighborhoods в FILO2 \source{8} --- принцип, применимый к крупномасштабному mTSP}
\end{figure}

\subsection{ITSHA и двухэтапные mTSP-эвристики из литературы}

Статья Zheng и соавторов посвящена итеративному двухэтапному эвристическому алгоритму для mTSP \source{10}. Для данной работы она особенно важна, потому что рассматривает именно multiple traveling salesmen problem и обсуждает как MINSUM, так и MINMAX-ориентированные варианты. Общая логика таких подходов состоит в разделении задачи на две части: сначала нужно получить разумное распределение вершин между коммивояжерами, затем улучшить порядок обхода и само распределение локальными ходами.

В ITSHA начальное решение строится с учетом кластерной структуры данных. Это близко к прикладной интуиции: если точки естественно образуют группы, то маршруты также должны использовать эту геометрию, а не распределять вершины случайно. После построения начального решения применяется вариативный локальный поиск, который меняет как порядок вершин в маршрутах, так и распределение между маршрутами. Важная идея статьи --- не ограничиваться одной простой процедурой улучшения, а чередовать несколько типов окрестностей, чтобы выходить из локальных минимумов.

По отношению к данному проекту ITSHA можно рассматривать как научное обоснование двух решений. Первое --- добавление cluster-aware начальных маршрутов в поздние версии \texttt{lkh-wrapper}. Второе --- необходимость межмаршрутных улучшений: только внутримаршрутного 2-opt недостаточно, если исходное назначение вершин агентам оказалось слабым. Поэтому развитие \texttt{v12} логично вести в сторону более богатого набора relocate/swap/block moves и контроля их стоимости через candidate sets.

\begin{figure}[H]
\centering
\fbox{\begin{minipage}{0.92\textwidth}
\centering
\textbf{Этап 1.} Построение начального решения с учетом кластерной структуры входных точек\\[4pt]
$\downarrow$\\[4pt]
\textbf{Этап 2.} Итеративное улучшение: чередование внутримаршрутных и межмаршрутных ходов (relocate, swap, 2-opt), выход из локальных минимумов через разнообразие окрестностей\\[4pt]
$\downarrow$\\[4pt]
Лучшее найденное решение MINSUM/MINMAX
\end{minipage}}
\caption{Двухэтапная структура ITSHA \source{10} и аналогия с \texttt{lkh-wrapper-v12}: cluster-aware начало и межмаршрутные переносы}
\end{figure}

\subsection{Кластерно-генетические подходы и MMCTSP}

Работа Bao и соавторов рассматривает min-max clustered traveling salesmen problem и предлагает двухэтапный генетический подход \source{11}. Эта постановка отличается от текущей MINSUM-версии проекта, но полезна для анализа случаев, когда входные точки имеют выраженную кластерную структуру. В таких данных важно не только построить короткие маршруты, но и решить, какие кластеры должен обслуживать каждый агент.

Двухэтапная идея выглядит естественно. На первом уровне оптимизируется порядок обхода внутри кластеров или между кластерами, на втором --- распределение кластеров между коммивояжерами. Генетический алгоритм используется для поиска хороших комбинаций, поскольку прямой перебор всех назначений быстро становится невозможным. Для данной работы это направление важно как возможное расширение: если в генераторе инстансов явно присутствуют кластеры, то алгоритм может сначала работать на уровне групп точек, а затем уточнять маршруты внутри каждой группы.

В текущем проекте полноценный генетический алгоритм не реализован, однако некоторые идеи уже близки к использованным решениям. Например, cluster-aware построение начальных маршрутов в \texttt{lkh-wrapper-v12} играет похожую роль: оно пытается не разрушать естественную геометрию данных. В дальнейшем это можно усилить, добавив отдельный этап назначения кластеров агентам и сравнив его с чисто жадным построением.

\begin{figure}[H]
\centering
\fbox{\begin{minipage}{0.92\textwidth}
\centering
входные точки с кластерной структурой\\[4pt]
$\downarrow$ \textbf{Уровень кластеров}: оптимизация распределения кластеров по агентам (генетический оператор скрещивания и мутации)\\[4pt]
$\downarrow$ \textbf{Уровень вершин}: оптимизация порядка обхода внутри каждого кластера\\[4pt]
MINMAX-оптимизированное mTSP-решение
\end{minipage}}
\caption{Двухуровневая структура генетического подхода для clustered MMCTSP \source{11}; постановка MINMAX/clustered отличается от MINSUM в данной работе, но принцип иерархической декомпозиции применим}
\end{figure}

\subsection{POPMUSIC и построение candidate sets}

POPMUSIC-подход Taillard и Helsgaun важен для данной работы не как отдельный mTSP-решатель, а как аргумент в пользу качественного построения candidate sets \source{7}. На больших TSP-инстансах полный граф слишком дорог, поэтому качество локального поиска во многом зависит от того, какие ребра попадут в набор перспективных кандидатов. Если candidate set слишком узкий, алгоритм пропустит хорошие ходы. Если он слишком широкий, время работы снова станет неприемлемым.

В статье по POPMUSIC показано, что candidate edges можно строить на основе анализа локальных подзадач и объединения информации из нескольких туров. Для проекта это подтверждает общий инженерный выбор: ускорение должно достигаться не только сокращением числа итераций, но и более умным выбором ребер, которые алгоритм вообще рассматривает. Поэтому геометрические candidate sets, сеточное разбиение пространства и кэш расстояний в \texttt{lkh-wrapper} являются не временными техническими упрощениями, а центральной частью алгоритмической идеи.

\begin{figure}[H]
\centering
\fbox{\begin{minipage}{0.92\textwidth}
\centering
крупный граф с $n$ вершинами\\[4pt]
$\downarrow$ разбить на перекрывающиеся локальные подзадачи, оптимизировать каждую\\[4pt]
$\downarrow$ из каждой локальной задачи извлечь candidate edges (ребра, часто встречающиеся в хороших частичных решениях)\\[4pt]
объединить candidate sets для всего графа\\[4pt]
\textit{В проекте}: более простая геометрическая версия --- ближайшие $k$ соседей по координатам через пространственную сетку
\end{minipage}}
\caption{Принцип POPMUSIC \source{7}: построение candidate sets через локальные подзадачи; в \texttt{lkh-wrapper} применяется геометрическое приближение той же идеи}
\end{figure}

\subsection{Сравнительная характеристика эвристик}

Сводное сравнение методов показывает, что ни одна эвристика не является универсально лучшей во всех режимах. Простые методы важны для baseline и быстрой проверки. GRASP дает лучшее качество за счет многократных запусков и межмаршрутных обменов. LKH-3 задает сильный внешний ориентир качества. Собственные LKH-подобные версии занимают промежуточное место: они быстрее и лучше контролируются в экспериментальном контуре, но пока уступают полноценному LKH-3 по глубине поиска.

{\small
\begin{longtable}{>{\raggedright\arraybackslash}p{0.21\textwidth}>{\raggedright\arraybackslash}p{0.25\textwidth}>{\raggedright\arraybackslash}p{0.22\textwidth}>{\raggedright\arraybackslash}p{0.22\textwidth}}
\toprule
\textbf{Метод} & \textbf{Основная идея} & \textbf{Сильная сторона} & \textbf{Ограничение}\\
\midrule
\endhead
\texttt{rand+nn} & случайное разбиение и ближайший сосед внутри группы & минимальное время, простой baseline & слабое назначение вершин маршрутам\\
\texttt{2opt+greed} & жадное построение и 2-opt внутри маршрутов & быстро улучшает порядок обхода & почти не исправляет межмаршрутное распределение\\
\texttt{grasp} & RCL, много стартов, 2-opt и обмены & лучше выходит из локальных ошибок построения & дороже по времени и чувствителен к параметрам\\
LKH-3 & локальный поиск переменной глубины, candidate sets, penalty functions & сильный внешний ориентир качества & сложнее интеграция и меньше прозрачности для модификаций\\
\texttt{lkh-wrapper} & собственные candidate-based эвристики для mTSP & управляемость, скорость, time-budget & пока меньше глубина поиска, чем у LKH-3\\
FILO2 & гранулярные окрестности и кэш затронутых вершин & идеи масштабирования до очень больших данных & не является прямым mTSP-решателем проекта\\
ITSHA & двухэтапное построение и вариативный локальный поиск & хорошо соответствует mTSP и кластерной геометрии & требует более сложной реализации окрестностей\\
ALNS \source{24,25} & destroy/repair-операторы с адаптивными весами и SA-приёмом & гибкая архитектура, единый каркас для многих VRP-вариантов & не реализован в проекте; ориентир для развития\\
SISR \source{28} & ruin-and-recreate через удаление пространственных строк клиентов & SOTA на CVRP/VRPTW при минимальной сложности & требует concept'а string removal, не сводится к 2-opt\\
HGS \source{26,27} & популяционный поиск + Split-разделители + intensive education & SOTA для CVRP, эталон в DIMACS & архитектурно сложнее ALNS, требует управления популяцией\\
KGLS \source{30} & granular GLS на n до~30~000 & масштабируется до десятков тысяч клиентов & ориентирован на CVRP, не на чистый mTSP\\
Memetic search \source{31} & memetic-алгоритм для MINMAX-mTSP с EAX-кроссовером & SOTA для MINMAX-mTSP с одним и несколькими депо & MINMAX-постановка отличается от текущей\\
Кластерно-генетический подход & назначение кластеров и эволюционный поиск & полезен для явно кластерных данных & относится к отличающейся MINMAX/clustered-постановке\\
\bottomrule
\end{longtable}
}

Из этого сравнения следует направление дальнейшего улучшения проекта. Базовые алгоритмы уже выполняют свою роль как контрольные ориентиры. Основной потенциал развития находится в LKH-подобной линии: нужно усилить межмаршрутные ходы, аккуратно использовать кластерную структуру входных данных и перенести из FILO2/POPMUSIC идею более умного ограничения области поиска. Тогда можно приблизиться к качеству LKH-3, сохранив при этом управляемость собственного решателя и удобство массовых экспериментов. Как отдельный вектор развития для крупных инстансов целесообразно рассматривать интеграцию SISR-оператора \source{28} (string-removal как естественный destroy для перебалансировки маршрутов между агентами) и granular neighborhoods Toth--Vigo \source{29} как формальной замены текущим геометрическим candidate sets. Для приближения к SOTA на средних инстансах стоит изучить применимость HGS-каркаса \source{26,27} с адаптацией Split-алгоритма к mTSP-постановке.

\section{ПОСТАНОВКА ЗАДАЧИ И РЕАЛИЗАЦИЯ}

\subsection{Формальная постановка в проекте}

Входной файл mTSP-инстанса содержит число вершин $n$, число агентов $m$ и координаты всех вершин. Вершина с индексом $0$ считается депо. Требуется построить $m$ маршрутов, удовлетворяющих условиям:
\begin{enumerate}
  \item каждый маршрут начинается в депо и заканчивается в депо;
  \item каждая клиентская вершина $1,\ldots,n-1$ встречается ровно в одном маршруте;
  \item значение MINSUM минимизируется или, для приближенного алгоритма, становится как можно меньше при заданном бюджете времени.
\end{enumerate}

Валидность решения проверяется отдельно от вычисления целевой функции. Это принципиально важно: эвристика может вернуть численно короткие маршруты, но если какая-либо вершина пропущена или посещена дважды, такое решение нельзя учитывать в сравнении.

\subsection{Экспериментальный контур}

Репозиторий проекта содержит C++-реализации решателей и Python-скрипты для генерации данных, пакетного запуска и построения отчетов. Основные элементы контура:
\begin{itemize}
  \item \texttt{run\_mtsp.py} --- единая точка запуска mTSP-решателей;
  \item \texttt{experiments/generate\_mtsp\_instances.py} --- генерация синтетических инстансов;
  \item \texttt{experiments/run\_benchmarks.py} --- пакетный запуск алгоритмов по конфигурации;
  \item \texttt{python/mtsp\_baselines/ortools\_routing.py} --- внешний baseline на OR-Tools;
  \item \texttt{python/mtsp\_baselines/tsp\_transform\_lkh.py} --- baseline через построение TSP-тура и последующее разбиение на маршруты;
  \item \texttt{data/results/} --- CSV и JSON-артефакты экспериментов.
\end{itemize}

\begin{figure}[H]
\centering
\fbox{\begin{minipage}{0.92\textwidth}
\centering
Генерация инстансов $\rightarrow$ запуск решателей $\rightarrow$ проверка валидности маршрутов $\rightarrow$ пересчет MINSUM $\rightarrow$ агрегация CSV/JSON $\rightarrow$ сравнительный анализ
\end{minipage}}
\caption{Схема экспериментального контура}
\end{figure}

Такое разделение уменьшает риск смешения кода алгоритма и кода оценки. Решатель отвечает за построение маршрутов, а отдельная часть проекта проверяет корректность результата и пересчитывает objective по сохраненным маршрутам.

\subsection{Набор реализованных алгоритмов}

\begin{longtable}{>{\raggedright\arraybackslash}p{0.24\textwidth}>{\raggedright\arraybackslash}p{0.68\textwidth}}
\toprule
\textbf{Алгоритм} & \textbf{Роль в исследовании}\\
\midrule
\endhead
\texttt{rand+nn} & Быстрый конструктивный baseline: случайность и ближайший сосед. Нужен как нижняя граница качества для простых решений.\\
\texttt{2opt+greed} & Жадное построение с локальным 2-opt. Позволяет оценить вклад внутримаршрутного локального поиска.\\
\texttt{grasp} & Рандомизированный многостартовый подход с restricted candidate list, 2-opt и межмаршрутными обменами.\\
\texttt{lkh-wrapper-v1...v3} & Ранние LKH-подобные версии; используются для проверки эффекта перехода от простой обертки к candidate-based логике.\\
\texttt{lkh-wrapper-v4...v7} & Версии для крупных инстансов с геометрическими кандидатами, кэшированием расстояний, бюджетами времени и ускоренными процедурами улучшения.\\
\texttt{lkh-wrapper-v12} & MINSUM-ориентированная версия: строит быстрое первое решение, использует savings/cluster-aware идеи и принимает межмаршрутные изменения по суммарной длине маршрутов.\\
\texttt{OR-Tools GLS} & Внешний эвристический ориентир на малых и средних инстансах.\\
\texttt{TSP-transform + LKH} & Внешний baseline: строит TSP-тур по клиентам и затем разбивает его на маршруты динамическим программированием.\\
\texttt{LKH-3 baseline} & Наиболее сильный внешний ориентир для части ручных сравнений.\\
\bottomrule
\end{longtable}

\subsection{Особенности LKH-подобной реализации}

Поздние версии \texttt{lkh-wrapper} не являются прямым запуском LKH-3 внутри C++-кода. Это собственные эвристические обертки, вдохновленные LKH-подходом и адаптированные под быстрый экспериментальный контур. Их общая логика состоит из нескольких этапов:

\begin{enumerate}
  \item построение геометрических candidate sets: для малых инстансов возможен точный перебор ближайших соседей, для крупных применяется пространственная сетка;
  \item построение начальных маршрутов: используется жадная или кластерно-ориентированная процедура, учитывающая расстояние до депо и распределение точек;
  \item локальное улучшение маршрутов: 2-opt ограничивается кандидатными ребрами, чтобы не выполнять полный квадратичный просмотр;
  \item межмаршрутные изменения: вершины или блоки переносятся между маршрутами, если это уменьшает MINSUM или улучшает баланс;
  \item соблюдение бюджета времени: крупные циклы регулярно проверяют deadline и возвращают лучшее найденное решение.
\end{enumerate}

В \texttt{lkh-wrapper-v12} отдельно выделяется бюджет на получение первого решения и бюджет на улучшение. Такая схема важна для прикладного режима: даже если улучшение не успело завершиться, алгоритм должен вернуть валидное решение. В коде также используется кэш расстояний для крупных инстансов, так как хранение полной матрицы расстояний становится чрезмерным по памяти.

\begin{figure}[H]
\centering
\fbox{\begin{minipage}{0.92\textwidth}
\centering
Candidate sets $\rightarrow$ fast seed routes $\rightarrow$ savings/cluster seed $\rightarrow$ route sanitizing $\rightarrow$ local ILS $\rightarrow$ inter-route relocate/swap $\rightarrow$ final validation
\end{minipage}}
\caption{Упрощенная логика работы поздней LKH-подобной версии}
\end{figure}

% TODO (опционально): добавить PNG-визуализации маршрутов (uniform, clustered, outliers) из сохраненных JSON в data/results/.
% Скрипт визуализации: matplotlib + json.load(routes_file) -> scatter + LineCollection по маршрутам.
% Вставить как \includegraphics в конец раздела или в приложение.

\section{МЕТОДИКА ВЫЧИСЛИТЕЛЬНОГО ЭКСПЕРИМЕНТА}

\subsection{Семейства тестовых инстансов}

Чтобы не проверять алгоритмы только на одном типе геометрии, в проекте используются несколько семейств синтетических данных. Такой дизайн лучше отражает разные сценарии маршрутизации и снижает риск, что алгоритм случайно хорошо работает только на равномерных точках.

\small
\begin{longtable}{>{\raggedright\arraybackslash}p{0.32\textwidth}>{\raggedright\arraybackslash}p{0.53\textwidth}>{\raggedright\arraybackslash}p{0.08\textwidth}}
\toprule
\textbf{Семейство} & \textbf{Смысл проверки} & \textbf{Депо}\\
\midrule
\endhead
\texttt{uniform} & Равномерное распределение точек. Базовый сценарий без выраженной структуры. & центр\\
\texttt{clustered-center} & Несколько естественных групп клиентов. Проверяется умение алгоритма сохранять компактные группы. & центр\\
\texttt{clustered-offset-depot} & Кластерная структура при смещенном депо. Сценарий сложнее, потому что ближайшие кластеры и баланс маршрутов конфликтуют. & край\\
\texttt{mixed-outliers} / \texttt{mixed} & Основная масса точек находится в группах, но есть удаленные выбросы. Проверяется устойчивость к длинным хвостам. & центр\\
\texttt{outliers} & Усиленный сценарий с удаленными точками. Важен для анализа межмаршрутного перераспределения. & центр\\
\texttt{high-m-stress} & Повышенная нагрузка по числу агентов и маршрутов. Проверяется устойчивость распределения вершин между маршрутами. & центр\\
\bottomrule
\end{longtable}
\normalsize

В ранних экспериментах использовалась сетка $n \in \{20,50,100,200\}$, $m \in \{2,3,5\}$ и 10 повторов. Позднее она была расширена семействами инстансов и отдельными крупными задачами до $n=100000$.

\subsection{Метрики}

Основная метрика качества --- значение MINSUM. Для сравнения с внешним ориентиром используется относительное отклонение:
\begin{equation}
  \mathrm{gap}(A,B) = 100 \cdot \frac{F(A)-F(B)}{F(B)},
\end{equation}
где $F(A)$ --- MINSUM исследуемого алгоритма, а $F(B)$ --- MINSUM базового или reference-решения. Если gap положителен, исследуемый алгоритм хуже ориентира по MINSUM; если отрицателен, он лучше.

Дополнительно учитываются:
\begin{itemize}
  \item среднее время работы;
  \item число валидных запусков;
  \item число запусков, где алгоритм превзошел reference;
  \item баланс маршрутов, измеряемый как отношение максимальной длины маршрута к средней.
\end{itemize}

В работе важно не смешивать разные типы сравнения. OR-Tools с лимитом 5 секунд не является точным оптимумом. Это внешний эвристический ориентир. Поэтому выводы формулируются не как «алгоритм дал оптимальное решение», а как «алгоритм имеет такое-то отклонение от выбранного reference при таких-то условиях».

\subsection{Валидация результатов}

Каждый результат должен удовлетворять трем условиям: все маршруты начинаются и заканчиваются в депо, все клиентские вершины посещены ровно один раз, значение MINSUM пересчитано независимой функцией по сохраненным маршрутам. Такой подход особенно важен при сравнении нескольких поколений алгоритмов: ускоренная версия может выглядеть привлекательной по времени, но если она возвращает невалидные маршруты, ее нельзя учитывать в агрегатах качества.

Последовательность воспроизведения экспериментов состоит из трёх шагов. Сначала выполняется сборка проекта:
\begin{verbatim}
cmake -B build -S . -DCMAKE_BUILD_TYPE=Release
cmake --build build --config Release
\end{verbatim}
Затем запускается benchmark через Python-скрипт \texttt{python/mtsp\_runner.py} с нужными параметрами (инстанс, солвер, бюджет времени). Каждый запуск принимает параметры \texttt{--input-file}, \texttt{--step} (имя алгоритма), \texttt{--time-budget-ms} и \texttt{--seed}. Результаты сохраняются в CSV и JSON в папке \texttt{data/results}. Агрегированные отчеты формируются повторным запуском runner'а с флагом генерации сводки. Все инстансы для экспериментов хранятся в \texttt{data/mtsp} и воспроизводимы через генератор с фиксированным seed.

\section{РЕЗУЛЬТАТЫ И ИХ АНАЛИЗ}

\subsection{Базовое сравнение на равномерных инстансах}

Первый набор экспериментов сравнивал \texttt{rand+nn}, \texttt{2opt+greed}, \texttt{grasp} и \texttt{lkh-wrapper-v2} на равномерных инстансах малой и средней размерности. По агрегированным результатам \texttt{lkh-wrapper-v2} дал лучшее среднее значение MINSUM среди внутренних алгоритмов, при этом оставался значительно быстрее GRASP.

\begin{table}[H]
\centering
\caption{Агрегированное сравнение внутренних алгоритмов на равномерных инстансах}
\begin{tabular}{lrrrr}
\toprule
Алгоритм & Запуски & Средний MINSUM & Время, с & Валидные\\
\midrule
\texttt{rand+nn} & 10 & 1571.74 & 0.0005 & 10\\
\texttt{2opt+greed} & 10 & 1042.94 & 0.0007 & 10\\
\texttt{grasp} & 10 & 866.91 & 1.0469 & 10\\
\texttt{lkh-wrapper-v2} & 10 & 821.33 & 0.0181 & 10\\
\bottomrule
\end{tabular}
\end{table}

Результат подтверждает ожидаемую иерархию методов. Чистый randomized nearest neighbor работает почти мгновенно, но имеет худшее качество. Добавление 2-opt резко улучшает результат. GRASP дает еще более качественные маршруты, но платит за это временем. \texttt{lkh-wrapper-v2} оказывается наиболее удачным компромиссом в этом наборе: его средний MINSUM ниже, чем у GRASP, а среднее время меньше примерно на два порядка.

\subsection{Сравнение с OR-Tools}

Для внешнего ориентира использовался OR-Tools Routing с Guided Local Search и лимитом 5 секунд. В таком сравнении все внутренние методы в среднем уступают OR-Tools по MINSUM, но разрыв существенно различается.

\begin{table}[H]
\centering
\caption{Средний gap к OR-Tools на равномерных инстансах}
\small
\begin{tabular}{lrrr}
\toprule
Алгоритм & Gap, \% & Время, с & Лучше ref.\\
\midrule
\texttt{rand+nn} & 114.73 & 0.0005 & 0.00\\
\texttt{2opt+greed} & 47.41 & 0.0007 & 0.00\\
\texttt{grasp} & 20.77 & 1.0469 & 0.00\\
\texttt{lkh-wrapper-v2} & 14.48 & 0.0181 & 0.08\\
\bottomrule
\end{tabular}
\normalsize
\end{table}

Эти данные не означают, что OR-Tools всегда лучше при любом временном бюджете. Его лимит в данном сравнении составляет 5 секунд, тогда как часть внутренних алгоритмов работает за миллисекунды. Однако таблица полезна для оценки качества: она показывает, насколько далеко находятся быстрые эвристики от сильного внешнего решения.

\subsection{Расширенный benchmark по семействам инстансов}

После добавления нескольких семейств инстансов картина стала более содержательной. На расширенном наборе \texttt{lkh-wrapper-v2} и \texttt{lkh-wrapper-v3} заметно превосходят ранние версии, GRASP и простые baseline-алгоритмы по среднему MINSUM. При этом \texttt{lkh-wrapper-v3} обычно быстрее \texttt{v2}, но по качеству немного уступает или находится рядом.

\begin{table}[H]
\centering
\caption{Средние значения по расширенному multi-family benchmark}
\begin{tabular}{lrr}
\toprule
Алгоритм & Средний MINSUM & Среднее время, с\\
\midrule
\texttt{rand+nn} & 1456.23 & 0.0007\\
\texttt{2opt+greed} & 999.30 & 0.0008\\
\texttt{grasp} & 820.59 & 1.5138\\
\texttt{lkh-wrapper-v1} & 875.99 & 0.0395\\
\texttt{lkh-wrapper-v2} & 677.61 & 0.0198\\
\texttt{lkh-wrapper-v3} & 683.44 & 0.0149\\
\bottomrule
\end{tabular}
\end{table}

В разрезе семейств видно, что преимущество LKH-подобных версий сохраняется не только на равномерных точках. Особенно сильный эффект наблюдается на кластерных данных: \texttt{lkh-wrapper-v2} снижает средний MINSUM относительно GRASP и ранних LKH-оберток, что согласуется с идеей использования геометрической структуры и более направленного локального поиска.

\begin{table}[H]
\centering
\caption{Средний MINSUM по семействам для ключевых алгоритмов}
\begin{tabular}{lrrrr}
\toprule
Семейство & \texttt{2opt+greed} & \texttt{grasp} & \texttt{lkh-v2} & \texttt{lkh-v3}\\
\midrule
\texttt{clustered-center} & 811.99 & 663.14 & 536.62 & 539.58\\
\texttt{clustered-offset-depot} & 959.50 & 817.75 & 538.31 & 550.34\\
\texttt{mixed-outliers} & 1072.43 & 863.14 & 765.89 & 771.07\\
\texttt{uniform} & 1153.26 & 938.34 & 869.62 & 872.79\\
\bottomrule
\end{tabular}
\end{table}

\subsection{Эволюция LKH-подобных версий}

Сравнение \texttt{lkh-wrapper-v1}, \texttt{v2} и \texttt{v3} показывает, что переход от начальной обертки к версиям с более аккуратной candidate-based логикой дал существенный прирост качества. Средний gap к OR-Tools снизился примерно с 29.63\% у \texttt{v1} до 14.48\% у \texttt{v2}. Версия \texttt{v3} почти повторила качество \texttt{v2}, но работала быстрее.

\begin{table}[H]
\centering
\caption{Сравнение ранних LKH-подобных версий с OR-Tools}
\small
\begin{tabular}{lrrr}
\toprule
Версия & Gap к OR-Tools, \% & Время, с & Валидные\\
\midrule
\texttt{lkh-wrapper-v1} & 29.63 & 0.0297 & 10\\
\texttt{lkh-wrapper-v2} & 14.48 & 0.0186 & 10\\
\texttt{lkh-wrapper-v3} & 14.61 & 0.0102 & 10\\
\bottomrule
\end{tabular}
\normalsize
\end{table}

Это важный результат для курсовой работы: он показывает не просто наличие нескольких запусков, а инженерную динамику. В проекте есть измеримый эффект от изменения алгоритма, а не только разовое сравнение готовых библиотек.

\subsection{Крупные инстансы и проблема масштабирования}

На крупных инстансах размером 10000--100000 вершин основная сложность смещается от качества локальных перестановок к способности алгоритма вообще успеть построить валидное решение и улучшить его в пределах бюджета. Эксперименты с версиями \texttt{v4}, \texttt{v5}, \texttt{v6} показывают, что ускорение может сопровождаться потерей качества или валидности.

\begin{table}[H]
\centering
\caption{Версии v4--v6 на крупных uniform/mixed инстансах с бюджетом около 30 секунд}
\small
\begin{tabular}{lrrr}
\toprule
Версия & MINSUM по валидным & Время, с & Валидность\\
\midrule
\texttt{lkh-wrapper-v4} & $2.38\cdot 10^8$ & 42.87 & 1.00\\
\texttt{lkh-wrapper-v5} & $2.63\cdot 10^8$ & 30.09 & 1.00\\
\texttt{lkh-wrapper-v6} & $3.37\cdot 10^5$ & 28.34 & 0.25\\
\bottomrule
\end{tabular}
\normalsize
\end{table}

Среднее значение \texttt{v6} в этой таблице нельзя трактовать как преимущество качества, потому что валидными оказались только отдельные запуски. Более корректный вывод заключается в другом: версия \texttt{v6} показала нестабильность, а \texttt{v4} и \texttt{v5} возвращали валидные решения, но \texttt{v5} обычно была быстрее ценой ухудшения MINSUM. Следовательно, для итогового алгоритма требуется одновременно контролировать валидность, time-budget и качество.

\subsection{Сравнение v12 с внешним LKH-3 baseline}

Поздняя версия \texttt{lkh-wrapper-v12} была ориентирована на MINSUM и быстрый возврат валидного решения. Ручные сравнения с \texttt{lkh3-baseline} показывают отчетливый компромисс: внешний LKH-3 дает более короткие маршруты, но собственная версия работает существенно быстрее и часто строит более сбалансированные маршруты.

\begin{table}[H]
\centering
\caption{Сравнение \texttt{lkh-wrapper-v12} с \texttt{lkh3-baseline} на n=5000, m=5}
\begin{tabular}{lrrrr}
\toprule
Инстанс & LKH-3 & v12 & v12, с & Ускорение v12\\
\midrule
\texttt{uniform} & 127094.46 & 159615.23 & 4.80 & 7.64\\
\texttt{mixed} & 62667.81 & 75288.82 & 4.83 & 5.81\\
\texttt{outliers} & 60255.93 & 71603.77 & 4.80 & 5.72\\
\texttt{clustered-center} & 27837.74 & 37140.17 & 4.73 & 5.98\\
\texttt{high-m-stress} & 61312.38 & 84752.47 & 4.83 & 5.86\\
\bottomrule
\end{tabular}
\end{table}

В среднем на этих пяти инстансах LKH-3 лучше по MINSUM примерно на 21\%, а \texttt{v12} быстрее примерно в 6.2 раза. Это не подтверждает сильную версию гипотезы о превосходстве собственного LKH-подобного алгоритма над LKH-3 по качеству, но подтверждает более практичный вывод: при жестком бюджете времени собственная версия может давать валидное решение быстро, а дальнейшая работа должна быть направлена на сокращение разрыва по objective.

\begin{table}[H]
\centering
\caption{Баланс маршрутов в сравнении LKH-3 и v12 на n=5000}
\begin{tabular}{lrr}
\toprule
Инстанс & Imbalance LKH-3 & Imbalance v12\\
\midrule
\texttt{uniform} & 4.99 & 1.01\\
\texttt{mixed} & 3.85 & 2.74\\
\texttt{outliers} & 3.88 & 2.64\\
\texttt{clustered-center} & 1.73 & 1.04\\
\texttt{high-m-stress} & 4.99 & 1.11\\
\bottomrule
\end{tabular}
\end{table}

Баланс маршрутов не является основной целью при MINSUM, однако таблица показывает интересный побочный эффект. \texttt{v12} чаще строит более равномерные маршруты, тогда как LKH-3 в режиме MINSUM может концентрировать больше нагрузки в отдельных маршрутах. Если прикладная задача потребует одновременно низкий MINSUM и ограничение на максимальную длину маршрута, эту особенность можно использовать как основу для следующей версии алгоритма.

Это наблюдение согласуется с известными теоретическими результатами о связи MINSUM и MINMAX-целей. Классическая работа Frederickson, Hecht и~Kim~\source{21} показывает, что разбиение оптимального TSP-тура на $m$~равных по длине частей даёт MINMAX-аппроксимацию с константным фактором, и при росте~$m$ относительный overhead на отдельный маршрут уменьшается. Carlsson и соавт.~\source{22} доказывают, что при равномерном распределении точек и большом~$m$ region-partition (балансированное разбиение области по нагрузке) даёт MINMAX-маршруты, чья суммарная длина асимптотически приближается к MINSUM-оптимуму. Эмпирический анализ Matl, Hartl и~Vidal~\source{23} на широком наборе VRP-инстансов показывает, что сокращение range длин маршрутов на~40\% достигается ценой увеличения суммарной стоимости лишь на~$\sim$2\%, а~37\% Pareto-эффективных решений лежат в пределах~10\% от cost-оптимума. С другой стороны, Bertazzi, Golden и~Wang~\source{20} в worst-case-анализе показывают, что отношение MINSUM-стоимости MINMAX-оптимума к MINSUM-оптимуму ограничено сверху $m$ и эта граница достигается на специально подобранных инстансах. Таким образом, наблюдаемое в \texttt{v12} существенное улучшение баланса при умеренном проигрыше по~MINSUM (порядка~21\%) находится в пределах теоретического коридора и подтверждает осмысленность разработки алгоритма, явно учитывающего MINMAX-структуру решения. Гипотеза о том, что при больших~$m$ MINMAX-ориентированная оптимизация будет давать MINSUM-результаты, близкие к прямой MINSUM-оптимизации, может быть проверена в дальнейших экспериментах с режимом~$m \geq 20$.

\subsection{Развитие архитектуры до v21: ALNS, Simulated Annealing и Parallel Tempering}

Версии \texttt{lkh-wrapper-v13}--\texttt{lkh-wrapper-v21} представляют собой не инкрементальное улучшение \texttt{v12}, а её архитектурное переоформление. Главным наблюдением, полученным при анализе \texttt{v12} на крупных инстансах (раздел~5.5), было то, что одного лишь \emph{улучшения} (intra-route 2-opt + межмаршрутный relocate) недостаточно: алгоритм быстро сходится к локальному минимуму и далее не имеет встроенного механизма ухода из него. Поэтому в \texttt{v21} была спроектирована модульная архитектура с явным разделением ответственности между построением, локальным поиском и метаэвристическим уровнем. Кодовая база разделена на 22~нумерованных модуля в \texttt{src/v21/core/} (\texttt{00\_types.hpp} -- \texttt{21\_popmusic.hpp}) и две отдельные точки входа --- \texttt{minsum/minsum\_solver.cpp} и \texttt{minmax/minmax\_solver.cpp}, регистрируемые в \texttt{SolverFactory} как \texttt{lkh\_v21\_minsum} и \texttt{lkh\_v21\_minmax}.

\textit{Конструктивная фаза.} Сначала строится несколько начальных решений разными методами (full round-robin nearest neighbor; быстрая версия по k-NN-кандидатам; polar sweep; balanced k-means для $n \leq 12{,}000$; savings-seed по Кларку--Райту для $n \leq 20{,}000$) и выбирается лучшее по \texttt{ScalarCost}. После этого выполняется параллельный final-2-opt по маршрутам и параллельный intra-route ILS (parallel polish) на OpenMP, число потоков задаётся autotune'ом.

\textit{Гранулярные кандидаты.} Кандидатные множества $C_v$ для каждой вершины строятся один раз на этапе \texttt{BuildKnnCandidates} через KDTree-2D; число соседей $k_{\mathrm{NN}}$ зависит от размера инстанса (от~26 для $n \leq 2000$ до~18 для $n > 60{,}000$). Все последующие move-операторы (relocate, swap, repair, intra-3opt, NeighborList2Opt) ограничивают рассматриваемые ходы только парами вершина--кандидат, что является непосредственной реализацией принципа Toth--Vigo~\source{29} для mTSP.

\textit{ALNS-каркас.} Главный цикл организован как Adaptive Large Neighborhood Search в стиле Ropke--Pisinger~\source{24,25}: на каждой итерации случайно выбирается пара (destroy-оператор, repair-оператор) с весом, адаптируемым по успехам в предыдущем сегменте. Реализованы четыре destroy-оператора (\texttt{Random}, \texttt{ClusterBfs}, \texttt{Expensive}, \texttt{Zone}) и два repair-оператора (\texttt{Cheapest}, \texttt{Regret2}). Размер разрушения $K_{\mathrm{destroy}}$ растёт с длиной streak'а no-improvement, что обеспечивает диверсификацию без необходимости полного рестарта. После каждого хода выполняется selective intra-route 2-opt только на маршрутах, помеченных как dirty, что амортизирует стоимость локального поиска.

\textit{Simulated Annealing с Ben-Ameur'ом.} Принятие/отклонение хода управляется отдельным модулем \texttt{SaEngine}, в котором температура $T_0$ оценивается на основе пилотного прогона из 100--200 случайных ходов и формулы $T_0 = -\bar{\Delta}^{+} / \log(p^{*})$, где $\bar{\Delta}^{+}$~--- средняя положительная дельта, $p^{*}$~--- целевая acceptance-вероятность \source{36}. Геометрическое охлаждение (cooling factor 0.99--0.997 в зависимости от tier'а) и встроенный reheat-механизм при затяжной стагнации обеспечивают баланс exploration/exploitation. \textit{Reheat} триггерится по двум условиям одновременно: накопление streak'а итераций без улучшения и порог wall-time без улучшения; на крупных инстансах ($n > 60{,}000$), где iter-rate низкий, время-основанный триггер становится основным механизмом escape'а из плато.

\textit{Parallel Tempering для среднего диапазона.} Для $12{,}000 < n \leq 60{,}000$ запускается ансамбль из 4~реплик с геометрической температурной лестницей $T_r = T_0 \cdot 8^{r/(K-1)}$; каждая реплика прогоняет ALNS-SA на собственном независимом бюджете, периодически (после каждой эпохи) делается попытка swap'а между соседними репликами по правилу Метрополиса $p_{\mathrm{swap}} = \min\{1, e^{(1/T_i - 1/T_j)(E_i - E_j)}\}$. На малых ($n \leq 12{,}000$) и очень больших ($n > 60{,}000$) инстансах PT отключён: эмпирически на таких размерностях затраты на синхронизацию реплик и копирование маршрутов перевешивают вклад от смешивания.

\textit{Дополнительные модули.} В архитектуру включены: GLS edge penalties (\texttt{19\_gls.hpp}) для штрафования часто посещаемых рёбер при стагнации; периодическая re-оптимизация пары наиболее длинных маршрутов (\texttt{20\_route\_pair\_reopt.hpp}); опциональная zone-decomposition через POPMUSIC (\texttt{21\_popmusic.hpp})~--- последний модуль реализован, но эмпирически отключён для $n > 60{,}000$, поскольку cheapest-insertion + bounded 2-opt на маршрутах длиной $\sim$20\,000 не превосходит уже достигнутый локальный оптимум. Все параметры (число реплик, $k_{\mathrm{NN}}$, размеры destroy, доли бюджета по фазам, пороги reheat) выбираются единым детерминированным dispatcher'ом \texttt{ResolveParamsForInstance}$(n, m, \mathrm{is\_minmax})$ с четырьмя tier'ами по~$n$.

\subsection{Variance audit и диагностика на flagship-инстансах}

Для строгой оценки стабильности \texttt{v21} был проведён formal variance audit на трёх flagship-инстансах из семейства \texttt{uniform} с $m=5$ и фиксированными временными бюджетами, соответствующими прикладным сценариям (среднее время на инстанс размером 10\,000~--- 60~секунд; 50\,000~--- 180~секунд; 100\,000~--- 380~секунд). Каждый инстанс прогонялся с 10~различными seed'ами; запуски проводились последовательно, чтобы исключить интерференцию по системным ресурсам. Общее wall-time эксперимента: 90.5~минут.

\begin{table}[H]
\centering
\caption{Variance audit \texttt{lkh\_v21\_minsum}: 10 seeds на flagship-инстанс}
\small
\begin{tabular}{lrrrrrr}
\toprule
Инстанс & Среднее & Std & cv, \% & Min & Max & Время, с\\
\midrule
\texttt{uniform\_n10000\_m5} & 410\,131.09 & 2\,438.79 & 0.59 & 407\,058.07 & 414\,022.71 & 32.6\\
\texttt{uniform\_n50000\_m5} & 4\,557\,677.51 & 7\,727.87 & 0.17 & 4\,546\,867.47 & 4\,567\,060.37 & 139.1\\
\texttt{uniform\_n100000\_m5} & 13\,009\,235.05 & 35\,955.67 & 0.28 & 12\,967\,085.51 & 13\,062\,844.89 & 370.9\\
\bottomrule
\end{tabular}
\normalsize
\end{table}

Относительное стандартное отклонение MINSUM (cv) находится в коридоре 0.17--0.59\%, что существенно ниже типичной величины эффекта от изменения архитектуры (несколько процентов на одно изменение). Это означает, что на данных инстансах \texttt{v21} обладает свойством, необходимым для строгого попарного сравнения версий: при бюджете 10~seeds любые улучшения порядка $\geq$0.5\% на $n \in \{50{,}000; 100{,}000\}$ должны быть статистически детектируемыми по парному критерию Wilcoxon signed-rank.

\textit{Анализ throughput-кривой при росте $n$.} Помимо итоговых значений MINSUM, дополнительной инструментацией собиралась трасса $(t_{\mathrm{ms}}, \mathrm{best\_cost})$ при каждом обновлении best-решения на протяжении ALNS-фазы. Сводные показатели по 10~seeds:

\begin{table}[H]
\centering
\caption{Диагностика анытайм-трасс по инстансам}
\small
\begin{tabular}{lrrrr}
\toprule
Инстанс & Iters/сек & best/iter, \% & sa\_reheats & Seeds с плато$^*$\\
\midrule
\texttt{uniform\_n10000\_m5} & 28.2 & 54--60 & 0/10 & 0/10\\
\texttt{uniform\_n50000\_m5} & 11.9 & 37--45 & 0/10 & 0/10\\
\texttt{uniform\_n100000\_m5} & 0.76 & 5--57 & 0/10 & 7/10\\
\bottomrule
\end{tabular}
\normalsize\\[2pt]
\footnotesize $^*$~Seed считается плато-кандидатом, если относительное падение MINSUM в последнем~20\% ALNS-фазы не превышает~0.1\%.
\end{table}

Здесь обнаруживаются три структурных особенности \texttt{v21}, существенные для интерпретации результатов и направлений дальнейшей работы.

\textit{Super-linear scaling.} При росте $n$ с $10{,}000$ до $100{,}000$ throughput падает в 21.6~раза, тогда как при гипотезе $O(n)$~per~iter ожидалось бы $\sim$10-кратное падение. Эмпирический показатель степени per-iter cost составляет $\approx n^{1.33}$, super-linear factor $\approx 2.14$. Этот overhead не локализован в одной функции и не объясняется простым багом в delta-evaluation: гранулярные кандидаты, инкрементальная длина маршрутов в \texttt{RouteList} и kNN-ограниченный 2-opt (\texttt{NeighborList2Opt}) уже корректны. Источник~--- внутренняя стоимость \texttt{NeighborList2Opt} на длинных маршрутах ($L \approx n/m$): принятые 2-opt-ходы выполняют \texttt{std::reverse} с амортизированной сложностью $O(L)$ на ход, что на маршрутах $L \approx 20{,}000$ доминирует над всеми остальными расходами. Это \emph{архитектурное} ограничение, требующее структурной переработки локального поиска (например, переход на doubly-linked list или splay-tree представление маршрута), а не точечной оптимизации.

\textit{PT отключён на $n > 60{,}000$.} На $n = 100{,}000$ во всех 10~seeds \texttt{pt\_replicas} = 1, что согласуется с настройкой autotune'а: при per-iter cost $> 1$~секунда амортизация на 4--6~реплик невыгодна, и одиночная реплика с большим бюджетом эмпирически даёт лучший результат. Это означает, что для крупных инстансов диверсификация должна обеспечиваться \emph{внутри} одной реплики (через destroy-операторы, GLS, reheat), а не через ансамблевый swap.

\textit{Реальная семантика плато.} Из 10~seeds на $n = 100{,}000$ \emph{семь} имеют относительное падение MINSUM в последнем~20\% ALNS-фазы менее~0.1\%, а у seed~6 период без улучшения составил 204~секунды из 242~секунд ALNS-фазы (84.5\%). При этом счётчик \texttt{sa\_reheats} равен нулю на всех 10~seeds: iter-count-based порог reheat'а ($120$~итераций без улучшения) на крупных инстансах не накапливается, поскольку при iter-rate $\sim$0.8~iter/сек 120~итераций без улучшения это $\approx$150~секунд, а к этому моменту бюджет уже близок к исчерпанию. Это мотивировало добавление в \texttt{v21} время-основанного триггера reheat'а как дополнительного условия (порог 30~секунд wall-time без улучшения для tier'а $n > 60{,}000$), но как самостоятельный архитектурный приём он не отменяет описанное super-linear-ограничение, лишь увеличивает шансы escape'а из плато для seeds с высокой долей plateau-времени.

Таким образом, variance audit подтверждает количественную стабильность \texttt{v21} на flagship-инстансах и одновременно идентифицирует точку, в которой дальнейший прогресс не сводится к настройке параметров и требует пересмотра представления маршрута на архитектурном уровне.

\subsection{FILO2-inspired стабилизация для high-m MINSUM}

Variance audit раздела~5.8 проводился на инстансах с малым числом агентов ($m=5$), где задача имеет много свободы в распределении вершин по маршрутам. Сценарий с большим числом агентов (high-m, $m \geq 80$) обнажил отдельную инженерную проблему \texttt{lkh\_v21\_minsum}: на инстансе \texttt{clustered-offset-depot\_n10000\_m100\_r01} solver систематически концентрировал 400--500 клиентов в одном маршруте, тогда как остальные оставались около ожидаемой нагрузки в 100~клиентов. Чисто математически постановка MINSUM может предпочитать такие несбалансированные конфигурации, если суммарная длина маршрутов при этом ниже; на практике подобный <<жирный>> маршрут добавляет нетривиальный intra-route detour cost и, как показали эксперименты, ухудшает MINSUM по сравнению с более равномерным распределением.

В качестве ориентира был использован FILO2~\source{8}, эталонный CVRP-решатель, чей адаптер для mTSP задаёт ёмкость маршрута как \texttt{capacity} = $\lceil (n-1)/m \rceil$. На том же инстансе FILO2 даёт MINSUM~17\,874 при максимальной длине маршрута~299 (imbalance~1.67), тогда как \texttt{lkh\_v21\_minsum}~--- 23\,261 при максимуме~473 (imbalance~2.04). Анализ исходного кода FILO2 (\texttt{solution/savings.hpp}, \texttt{localsearch/OneZeroExchange.hpp}, \texttt{localsearch/TailsExchange.hpp}) показал общий приём: ёмкость маршрута является \emph{предусловием} каждого move-оператора, проверяется перед расчётом дельты стоимости и стоит одно целочисленное сравнение.

Этот приём перенесён в \texttt{v21} как отдельный solver \texttt{lkh\_v21\_minsum\_cap}, без переписывания основной линии. Изменения локализованы в пяти файлах (новый файл-solver плюс минимальные правки в существующих): \texttt{05\_route\_list.hpp} (метод \texttt{RouteSize}), \texttt{14\_repair\_ops.hpp} (cap-skip в обоих repair-операторах + cap-aware fallback), \texttt{17\_pipeline.hpp} (проброс cap в \texttt{RepairContext} в обеих ветках~--- single-replica и PT), \texttt{18\_autotune.hpp} (поле \texttt{route\_cap}, default~0~--- legacy-режим), и новый \texttt{src/v21/minsum/minsum\_solver\_cap.cpp}. Существующий \texttt{lkh\_v21\_minsum} остаётся неизменным.

Cap вычисляется как $\lceil \mathrm{target} \cdot (1 + \mathrm{slack}) \rceil$, где $\mathrm{target} = \lceil (n-1)/m \rceil$, а \texttt{slack} настраивается через CLI-опцию \texttt{--route-cap-slack}. Эмпирический подбор на flagship-инстансе (одноseed-sweep, 60~с) показал нетривиальный оптимум:

\begin{table}[H]
\centering
\caption{Подбор slack на \texttt{clustered-offset-depot\_n10000\_m100} (seed=1, 60~с)}
\small
\begin{tabular}{lrrrr}
\toprule
slack & cap & MINSUM & max\_size & at-cap routes\\
\midrule
без cap (\texttt{v21}) & $\infty$ & 22\,158 & 513 & ---\\
0.10 & 111 & 21\,642 & 736 & 71\\
0.25 & 125 & 21\,188 & 463 & 25\\
0.50 & 150 & 20\,502 & 506 &  9\\
\textbf{0.75} & \textbf{175} & \textbf{20\,215} & \textbf{343} & \textbf{1}\\
1.00 & 200 & 20\,712 & 309 &  1\\
1.50 & 250 & 20\,974 & 322 &  0\\
\bottomrule
\end{tabular}
\normalsize
\end{table}

Слишком жёсткий cap (slack 0.10--0.25) ухудшает MINSUM, поскольку заставляет repair раз за разом отказываться от хороших candidate-positions и направлять клиентов в маршруты, географически удалённые от их естественных соседей. Slack=0.75 оказался <<точкой минимума>>: cap срабатывает крайне редко (в среднем 1~маршрут из 100 достигает порога), но самим фактом существования предотвращает образование жирного маршрута. Это согласуется с философией FILO2~--- ёмкость как hard precondition для редкого экстремума, а не как тесная квота.

С settled \texttt{slack=0.75} \texttt{lkh\_v21\_minsum\_cap} был сравнён с \texttt{lkh\_v21\_minsum} на четырёх high-m инстансах ($n=10\,000$, $m=100$, по 5~seeds на каждый, 60~с budget, оба запуска сегодня для контроля системного шума):

\begin{table}[H]
\centering
\caption{Парное сравнение \texttt{v21\_cap} vs \texttt{v21\_minsum} на high-m ($m=100$)}
\small
\begin{tabular}{lrrrrr}
\toprule
Инстанс & v21 mean & v21\_cap mean & $\Delta$ \% & seeds \checkmark & Wilcoxon $p$\\
\midrule
\texttt{clustered-offset-depot} & 21\,940 & 20\,677 & \textbf{$-5.76$} & 5/5 & \textbf{0.031}\\
\texttt{clustered-center} & 12\,972 & 12\,236 & \textbf{$-5.67$} & 5/5 & \textbf{0.031}\\
\texttt{n10000\_m100} (uniform) & 15\,669 & 15\,386 & $-1.81$ & 5/5 & \textbf{0.031}\\
\texttt{mixed-outliers} & 17\,414 & 16\,995 & $-2.40$ & 3/5 & 0.156\\
\bottomrule
\end{tabular}
\normalsize
\end{table}

Агрегированно по всем 20~парным запускам: средний $\Delta = -3.24\,\%$, 16~из~20~seeds улучшены, Wilcoxon signed-rank one-sided $p=0.002$ против гипотезы об отсутствии различия. Это устойчивый, статистически значимый эффект на $\alpha=0.05$.

Балансовые метрики ведут себя неоднородно по семействам. На двух инстансах из четырёх (flagship и \texttt{mixed-outliers}) cap не только улучшает MINSUM, но и резко уменьшает \texttt{max\_route} и \texttt{imbalance} (на \texttt{mixed-outliers} максимальный размер маршрута падает с 1116 до 478~клиентов, $-57\,\%$). На \texttt{clustered-center} и \texttt{n10000\_m100} баланс при этом немного ухудшается ($+5\,\%$ к \texttt{imbalance}): cap заполняет до cap некоторые <<геометрически удачные>> маршруты и направляет overflow в чуть более далёкие. Этот компромисс честно отражён в полном протоколе экспериментов (\texttt{docs/filo2\_inspired\_cap.md}); для отчёта важен сам факт того, что MINSUM-улучшение не достигается ценой взрыва \texttt{max\_route} как у LKH-3 (где \texttt{imbalance} достигает 59.5).

Разрыв до FILO2 на flagship-инстансе сократился с $\sim$30\,\% (\texttt{v21}: 23\,261 vs FILO2: 17\,874) до $\sim$15.7\,\% (\texttt{v21\_cap}: 20\,677). Это не паритет --- FILO2 продолжает сильно опережать на этих инстансах --- но более чем половина разрыва закрыта одним локальным изменением, не требующим архитектурной переработки.

\textit{Честная оговорка для постановки.} Чистая MINSUM-задача математически не имеет предпочтения к сбалансированным решениям; теоретически сильно несбалансированная конфигурация может иметь меньшую суммарную длину. Capacity-aware repair поэтому не является \emph{корректностной} правкой, а \emph{практической стабилизацией}: он убирает один класс патологических локальных оптимумов (<<жирный маршрут>>), которые ни один реальный диспетчер не примет даже если численная сумма ниже. Тот факт, что на каждом протестированном high-m инстансе это попутно снижает и саму MINSUM, является эмпирическим, а не теоретическим, результатом и согласуется с теоретическим коридором между MINSUM- и MINMAX-целями (раздел~5.6, \source{20,21,22,23}).

\subsection{Интерпретация результатов}

Выполненные эксперименты позволяют сформулировать несколько выводов.

Во-первых, простые эвристики полезны как baseline, но недостаточны для качественного решения mTSP. \texttt{rand+nn} и \texttt{2opt+greed} работают очень быстро, однако дают заметно худшую MINSUM.

Во-вторых, использование candidate-based локального поиска оправдано. Переход от ранних LKH-оберток к \texttt{v2}/\texttt{v3} почти вдвое снижает средний gap к OR-Tools на равномерных инстансах и дает устойчивый выигрыш на разных семействах данных.

В-третьих, масштабирование нельзя оценивать только по времени. Версия может быть быстрой, но невалидной или слишком сильно проигрывать по objective. Поэтому в отчете отдельно фиксируются valid runs и не делается вывод о качестве по строкам, где валидность низкая.

В-четвертых, внешний LKH-3 пока остается более сильным ориентиром по качеству. Собственная \texttt{v12} версия выигрывает по скорости и балансу, но проигрывает по MINSUM. Это честный и важный результат: он показывает, что проект находится не в состоянии «готовый промышленный решатель лучше всех», а в состоянии исследовательского прототипа с понятным направлением развития.

Для двух ключевых пар алгоритмов посчитано число инстансов, где первый алгоритм дал меньшее значение MINSUM (победа), большее (поражение) или одинаковое (ничья). Сравнение выполнено по идентичным инстансам из накопленных CSV-результатов.

\begin{table}[H]
\centering
\caption{Попарное сравнение алгоритмов: победы / поражения / ничьи по MINSUM}
\begin{tabular}{lrrr}
\toprule
Пара (алгоритм A vs алгоритм B) & Победы A & Поражения A & Ничьи\\
\midrule
\texttt{lkh-wrapper-v2} vs \texttt{grasp} & 93 & 27 & 0\\
\texttt{lkh-wrapper-v12} vs \texttt{lkh3-baseline} & 0 & 5 & 0\\
\bottomrule
\end{tabular}
\end{table}

Первая строка подтверждает преимущество \texttt{lkh-wrapper-v2} над \texttt{grasp} на большинстве инстансов~(93 из~120): использование candidate sets и LKH-подобного поиска стабильно дает лучший MINSUM, чем многостартовый жадный подход при одинаковых инстансах. Вторая строка описывает сравнение из раздела~5.3: на пяти рассмотренных инстансах $n=5000$ \texttt{lkh-wrapper-v12} уступает \texttt{lkh3-baseline} по качеству во всех случаях, что согласуется с наблюдаемым разрывом около~21\%.

\section{ДОСТОВЕРНОСТЬ, ОГРАНИЧЕНИЯ И РИСКИ}

Достоверность результатов обеспечивается несколькими мерами. Все алгоритмы запускаются через единый runner, что снижает риск различий в формате входа и выхода. Для каждого решения проводится независимая проверка посещения вершин и пересчет MINSUM. Результаты сохраняются в CSV/JSON, что позволяет повторно агрегировать данные без ручного переноса чисел. Сравнения выполняются на одних и тех же инстансах, поэтому различия между алгоритмами связаны с их поведением, а не с различием входных данных. Применённый дизайн~--- стратификация по семействам инстансов, парные сравнения на идентичных входных данных, независимая валидация маршрутов~--- соответствует рекомендациям современных стандартов экспериментального анализа алгоритмов \source{14,15,32,33}.

При этом работа имеет ограничения. Во-первых, часть экспериментов выполнена на синтетических данных. Это правильно для контролируемого исследования, но перед прикладным использованием алгоритмы нужно проверить на реальных географических наборах. Во-вторых, не для всех крупных запусков известна полная аппаратная конфигурация. В-третьих, некоторые сравнения имеют разный временной бюджет: OR-Tools запускался с лимитом 5 секунд, а собственные алгоритмы в ранних экспериментах часто работали значительно меньше. Поэтому выводы о «лучшем» алгоритме делаются отдельно для качества и отдельно для времени.

В-четвёртых, в текущей версии работы агрегированные таблицы не содержат стандартных отклонений и доверительных интервалов, а сравнение алгоритмов проводится через средние значения и парные счёты wins/losses без формальных статистических тестов. Это сознательное ограничение объёма работы: для строгого вывода о статистической значимости различий требуется применение непараметрических парных тестов (Wilcoxon signed-rank для пар алгоритмов, Friedman для множественного сравнения), что предусмотрено как направление дальнейшей работы. Размер выборки в~5 инстансов в сравнении \texttt{v12} с LKH-3 (раздел~5.6) также мал для уверенных выводов о среднем разрыве~21\%; разброс по инстансам составляет от $\sim$19\% до $\sim$38\%, и для устойчивой оценки требуется расширение набора. В-пятых, отсутствует anytime-анализ: представлены только итоговые значения MINSUM в момент истечения бюджета, без кривых time-vs-quality. Для итеративных эвристик на крупных инстансах ($n \geq 10^4$) такой анализ является более информативным критерием, чем точечная фиксация конечного значения~\source{34}.

Ещё одно ограничение связано с целью MINSUM. Она минимизирует суммарную длину маршрутов, но не гарантирует справедливого распределения нагрузки между агентами. В результатах видно, что LKH-3 может давать лучший MINSUM, но менее сбалансированные маршруты. Если практическая задача требует ограничить максимальную длину маршрута, постановку нужно расширить до MINMAX или многокритериальной цели.

Основной технический риск проекта --- рост времени локального поиска на больших инстансах. Для его снижения используются candidate sets, пространственная сетка, кэш расстояний и проверка бюджета времени. Второй риск --- невалидные решения в ускоренных версиях. Он снижается отдельной фазой sanitize/complete routes и независимой валидацией.

\section*{ЗАКЛЮЧЕНИЕ}
\addcontentsline{toc}{section}{ЗАКЛЮЧЕНИЕ}

В ходе работы была рассмотрена задача множественного коммивояжера с одним депо и целевой функцией MINSUM. Были изучены классические и современные подходы к TSP/mTSP, включая локальный поиск, GRASP, LKH, LKH-3, candidate sets и идеи масштабирования для крупных маршрутизационных задач. На основе этого был реализован экспериментальный контур, позволяющий генерировать инстансы, запускать несколько классов алгоритмов, проверять валидность маршрутов и агрегировать результаты.

Эксперименты показали, что LKH-подобные методы существенно превосходят простые baseline-алгоритмы по качеству решения. На равномерных инстансах \texttt{lkh-wrapper-v2} получил средний gap к OR-Tools около 14.48\%, тогда как GRASP имел около 20.77\%, \texttt{2opt+greed} --- 47.41\%, а \texttt{rand+nn} --- 114.73\%. На расширенном наборе семейств \texttt{lkh-wrapper-v2} и \texttt{v3} также дали лучшие средние значения MINSUM среди внутренних алгоритмов.

При переходе к крупным задачам была выявлена ключевая инженерная проблема: ускорение алгоритма должно контролироваться вместе с валидностью и качеством. Версии \texttt{v4}/\texttt{v5}/\texttt{v6} показали разные стороны этого компромисса: \texttt{v5} быстрее \texttt{v4}, но обычно хуже по MINSUM, а \texttt{v6} не обеспечила достаточную валидность. Поздняя MINSUM-ориентированная версия \texttt{v12} в ручных сравнениях с LKH-3 оказалась быстрее примерно в 6 раз на инстансах $n=5000$, но уступила по суммарной длине маршрутов примерно на 21\%. При этом \texttt{v12} часто давала более сбалансированные маршруты, что может быть полезно для дальнейшей многокритериальной постановки.

Архитектурный потолок \texttt{v12} (отсутствие встроенного механизма ухода из локального минимума и одиночная итеративная схема улучшения) мотивировал переход к существенно более сложной архитектуре в линейке \texttt{v13}--\texttt{v21}: 22~нумерованных модуля в \texttt{src/v21/core/} реализуют ALNS-каркас Ropke--Pisinger~\source{24,25} с четырьмя destroy-операторами и двумя repair-операторами, отдельный SaEngine с автонастройкой $T_0$ по Ben-Ameur'у~\source{36}, опциональный Parallel Tempering из 4~реплик для $n \in (12{,}000;\,60{,}000]$, GLS edge penalties и POPMUSIC-zone-decomposition. Variance audit на трёх flagship-инстансах ($n \in \{10{,}000; 50{,}000; 100{,}000\}$, по 10~seeds на каждый) показал относительное стандартное отклонение MINSUM 0.17--0.59\%~--- этот уровень шума достаточно низок, чтобы по парному критерию Wilcoxon signed-rank можно было детектировать улучшения порядка $\geq$0.5\%. Одновременно audit идентифицировал главное архитектурное ограничение: на $n = 100{,}000$ throughput ALNS-цикла составляет всего $\approx$0.8~итераций в секунду при super-linear scaling-факторе $\approx 2.14$ относительно $O(n)$, что приводит к плато на 7~из~10~seeds. Источник~--- стоимость \texttt{std::reverse} на длинных маршрутах внутри candidate-list 2-opt; точечной оптимизацией это не решается, требуется пересмотр представления маршрута (doubly-linked list или splay-tree).

В качестве отдельного направления для high-m сценариев ($m=100$) был реализован FILO2-вдохновлённый solver \texttt{lkh\_v21\_minsum\_cap} с capacity-aware repair (раздел~5.9). На четырёх инстансах $n=10\,000$, $m=100$ он даёт парный MINSUM-выигрыш в среднем 3.24\,\% (16/20 seeds улучшены, Wilcoxon $p=0.002$); на flagship-инстансе \texttt{clustered-offset-depot\_n10000\_m100} разрыв до FILO2 сократился с $\sim$30\,\% до $\sim$15.7\,\% при сохранении умеренного баланса маршрутов. Существующий \texttt{lkh\_v21\_minsum} остаётся неизменным; cap-вариант полезен как специализированный инструмент для high-m, когда главным риском является <<жирный маршрут>>.

Таким образом, цель работы в исследовательском смысле достигнута: разработан и проанализирован набор эвристических методов для mTSP, построен воспроизводимый экспериментальный контур, получены сравнения с базовыми и внешними алгоритмами, выявлены сильные и слабые стороны текущих решений. Рабочая гипотеза подтверждена частично. LKH-подобные эвристики действительно лучше простых конструктивных методов и GRASP на рассмотренных наборах, но пока не превосходят внешний LKH-3 по качеству MINSUM.

На основе проведённого исследования выделяются следующие направления дальнейшего развития проекта. \textbf{Алгоритмически}: интеграция SISR-оператора~\source{28} в качестве естественного destroy-механизма для перебалансировки маршрутов между агентами; формализация candidate sets через granular neighborhoods Toth--Vigo~\source{29} вместо текущей геометрической эвристики; адаптация HGS-каркаса~\source{26,27} с реализацией Split-алгоритма для mTSP-постановки; для крупных инстансов ($n \geq 10^4$)~--- внедрение KGLS-подхода~\source{30} с гранулярным GLS-поиском. \textbf{Методологически}: добавление парных непараметрических тестов (Wilcoxon, Friedman) и доверительных интервалов в агрегированные таблицы; построение anytime-кривых time-vs-quality для итеративных алгоритмов на крупных инстансах; расширение размера выборки в сравнении с LKH-3 с~5 до~30+ инстансов на каждый класс; добавление стандартных бенчмарков (например, Uchoa X-instances) для калибровки относительно SOTA. \textbf{По линии MINMAX}: проверка эмпирической гипотезы, что при больших $m$ ($\geq 20$) MINMAX-ориентированная оптимизация даёт MINSUM-результаты, близкие к прямой MINSUM-оптимизации, что согласовалось бы с теоретическими работами Frederickson--Hecht--Kim~\source{21}, Carlsson и соавт.~\source{22} и Matl--Hartl--Vidal~\source{23}. \textbf{По линии прикладного применения}: проверка алгоритмов на реальных географических наборах с дорожной метрикой и неравномерной плотностью клиентов, что особенно важно для оценки устойчивости candidate-based подходов в условиях, отличных от чистой евклидовой геометрии.

\newpage
\section*{СПИСОК ИСПОЛЬЗОВАННЫХ ИСТОЧНИКОВ}
\addcontentsline{toc}{section}{СПИСОК ИСПОЛЬЗОВАННЫХ ИСТОЧНИКОВ}

\begin{enumerate}
  \item Lin S., Kernighan B. W. An Effective Heuristic Algorithm for the Traveling-Salesman Problem // Operations Research. 1973. Vol. 21, No. 2. P. 498--516.
  \item Helsgaun K. An Effective Implementation of the Lin--Kernighan Traveling Salesman Heuristic // European Journal of Operational Research. 2000. Vol. 126, No. 1. P. 106--130.
  \item Helsgaun K. An Extension of the Lin--Kernighan--Helsgaun TSP Solver for Constrained Traveling Salesman and Vehicle Routing Problems. Roskilde University, 2017.
  \item Bektaş T. The Multiple Traveling Salesman Problem: An Overview of Formulations and Solution Procedures // Omega. 2006. Vol. 34, No. 3. P. 209--219.
  \item Jonker R., Volgenant T. Technical Note --- An Improved Transformation of the Symmetric Multiple Traveling Salesman Problem // Operations Research. 1988. Vol. 36, No. 1. P. 163--167.
  \item Croes G. A. A Method for Solving Traveling-Salesman Problems // Operations Research. 1958. Vol. 6, No. 6. P. 791--812.
  \item Taillard É. D., Helsgaun K. POPMUSIC for the Travelling Salesman Problem // European Journal of Operational Research. 2019. Vol. 272, No. 2. P. 420--429.
  \item Accorsi L., Vigo D. Routing One Million Customers in a Handful of Minutes. 2023. arXiv:2306.14205.
  \item Feo T. A., Resende M. G. C. Greedy Randomized Adaptive Search Procedures // Journal of Global Optimization. 1995. Vol. 6, No. 2. P. 109--133.
  \item Zheng J., Hong Y., Xu W., Li W., Chen Y. An Effective Iterated Two-stage Heuristic Algorithm for the Multiple Traveling Salesmen Problem. 2022. arXiv:2201.09424.
  \item Bao X., Wang G., Xu L., Wang Z. Solving the Min-Max Clustered Traveling Salesmen Problem Based on Genetic Algorithm // Biomimetics. 2023. Vol. 8, No. 238.
  \item Russell R. A. An Effective Heuristic for the M-Tour Traveling Salesman Problem with Some Side Conditions // Operations Research. 1977. Vol. 25, No. 3. P. 517--524.
  \item Reinelt G. TSPLIB --- A Traveling Salesman Problem Library // ORSA Journal on Computing. 1991. Vol. 3, No. 4. P. 376--384.
  \item Johnson D. S. A Theoretician's Guide to the Experimental Analysis of Algorithms // Data Structures, Near Neighbor Searches, and Methodology. 1999. P. 215--250.
  \item Hooker J. N. Testing Heuristics: We Have It All Wrong // Journal of Heuristics. 1995. Vol. 1, No. 1. P. 33--42.
  \item Google OR-Tools. Routing Library Documentation. Электронный ресурс: \url{https://developers.google.com/optimization/routing}.
  \item Dantzig G. B., Ramser J. H. The Truck Dispatching Problem // Management Science. 1959. Vol. 6, No. 1. P. 80--91.
  \item Toth P., Vigo D. Vehicle Routing: Problems, Methods, and Applications. SIAM, 2014.
  \item Материалы репозитория \texttt{mtsp-routing-optimization}: исходный код решателей, скрипты экспериментов и CSV/JSON-артефакты результатов. Локальный архив проекта, 2026.
  \item Bertazzi L., Golden B., Wang X. Min--Max vs. Min--Sum Vehicle Routing: A worst-case analysis // European Journal of Operational Research. 2015. Vol. 240, No. 2. P. 372--381.
  \item Frederickson G. N., Hecht M. S., Kim C. E. Approximation algorithms for some routing problems // SIAM Journal on Computing. 1978. Vol. 7, No. 2. P. 178--193.
  \item Carlsson J., Ge D., Subramaniam A., Wu A., Ye Y. Solving the min-max multi-depot vehicle routing problem // Lectures on Global Optimization (Pardalos~P., Coleman~T., eds.), Fields Institute Communications, vol.~55, AMS, 2009. P. 31--46.
  \item Matl P., Hartl R. F., Vidal T. Workload Equity in Vehicle Routing Problems: A Survey and Analysis // Transportation Science. 2018. Vol. 52, No. 2. P. 239--260.
  \item Ropke S., Pisinger D. An Adaptive Large Neighborhood Search Heuristic for the Pickup and Delivery Problem with Time Windows // Transportation Science. 2006. Vol. 40, No. 4. P. 455--472.
  \item Pisinger D., Ropke S. A general heuristic for vehicle routing problems // Computers \& Operations Research. 2007. Vol. 34, No. 8. P. 2403--2435.
  \item Vidal T., Crainic T. G., Gendreau M., Lahrichi N., Rei W. A Hybrid Genetic Algorithm for Multidepot and Periodic Vehicle Routing Problems // Operations Research. 2012. Vol. 60, No. 3. P. 611--624.
  \item Vidal T. Hybrid Genetic Search for the CVRP: Open-source implementation and SWAP* neighborhood // Computers \& Operations Research. 2022. Vol. 140, Article 105643.
  \item Christiaens J., Vanden Berghe G. Slack Induction by String Removals for Vehicle Routing Problems // Transportation Science. 2020. Vol. 54, No. 2. P. 417--433.
  \item Toth P., Vigo D. The Granular Tabu Search and Its Application to the Vehicle-Routing Problem // INFORMS Journal on Computing. 2003. Vol. 15, No. 4. P. 333--346.
  \item Arnold F., Gendreau M., Sörensen K. Efficiently solving very large-scale routing problems // Computers \& Operations Research. 2019. Vol. 107. P. 32--42.
  \item He P., Hao J.-K. Memetic search for the minmax multiple traveling salesman problem with single and multiple depots // European Journal of Operational Research. 2023. Vol. 307, No. 3. P. 1055--1070.
  \item Coffin M., Saltzman M. J. Statistical Analysis of Computational Tests of Algorithms and Heuristics // INFORMS Journal on Computing. 2000. Vol. 12, No. 1. P. 24--44.
  \item Bartz-Beielstein T., Chiarandini M., Paquete L., Preuss M. (Eds.). Experimental Methods for the Analysis of Optimization Algorithms. Springer, 2010.
  \item López-Ibáñez M., Stützle T. Automatically improving the anytime behaviour of optimisation algorithms // European Journal of Operational Research. 2014. Vol. 235, No. 3. P. 569--582.
  \item Cheikhrouhou O., Khoufi I. A comprehensive survey on the Multiple Traveling Salesman Problem: Applications, approaches and taxonomy // Computer Science Review. 2021. Vol. 40, Article 100369.
  \item Ben-Ameur W. Computing the Initial Temperature of Simulated Annealing // Computational Optimization and Applications. 2004. Vol. 29, No. 3. P. 369--385.
\end{enumerate}

\newpage
\section*{ПРИЛОЖЕНИЕ А. ЧЕКЛИСТ ФИНАЛЬНОЙ ПОДГОТОВКИ}
\addcontentsline{toc}{section}{ПРИЛОЖЕНИЕ А. ЧЕКЛИСТ ФИНАЛЬНОЙ ПОДГОТОВКИ}

Ниже перечислены элементы, которые остаётся проверить перед окончательной сдачей. Пункты, отмеченные $\checkmark$, уже включены в основной текст.

\begin{enumerate}
  \item[\checkmark] Аппаратная конфигурация экспериментов: Intel Core Ultra~7 165U, 32~ГБ RAM, Windows~11, Clang~20, Release --- добавлена в раздел <<Введение>>.
  \item[\checkmark] Текстовые схемы вместо рисунков из статей: LKH-3, FILO2, ITSHA, кластерный ГА, POPMUSIC --- добавлены в главу~2 в виде оформленных блоков \texttt{\textbackslash fbox}. При наличии PNG из \texttt{docs/materials} их можно заменить через \texttt{\textbackslash includegraphics}.
  \item[\checkmark] Таблица wins/losses/ties для пар \texttt{lkh-wrapper-v2} vs \texttt{grasp} и \texttt{v12} vs \texttt{lkh3-baseline} --- добавлена в раздел~5.4.
  \item[\checkmark] Команды воспроизведения экспериментов (сборка, запуск, сохранение результатов) --- добавлены в раздел~4.3.
  \item[\checkmark] Описание архитектуры \texttt{v13}--\texttt{v21}: ALNS-каркас с destroy/repair-операторами, SaEngine с автонастройкой $T_0$, Parallel Tempering на среднем tier'е, GLS, POPMUSIC --- добавлено в раздел~5.7.
  \item[\checkmark] Variance audit на flagship-инстансах ($n \in \{10\,000; 50\,000; 100\,000\}$, по 10 seeds): cv MINSUM 0.17--0.59\%, throughput-кривая, диагностика super-linear-scaling --- добавлено в раздел~5.8 со сводными таблицами.
  \item[\checkmark] Парный harness для повторяемых runs с резюме (\texttt{experiments/run\_audit.py}) и парный compare-script с Wilcoxon (\texttt{experiments/compare\_runs.py}) --- работают с готовыми run-JSON в \texttt{data/results/audit/baseline/}.
  \item[\checkmark] Anytime-инструментация: SaEngine эмитит трассу $(t_{\mathrm{ms}}, \mathrm{best\_cost})$ при каждом обновлении best, доступна как поле \texttt{anytime\_trace} в metadata.
  \item[\checkmark] FILO2-вдохновлённый \texttt{lkh\_v21\_minsum\_cap} solver: capacity-aware repair с $\mathrm{cap} = \lceil (n-1)/m \cdot 1.75 \rceil$. На 4~high-m инстансах ($m=100$): средний $\Delta = -3.24\,\%$ по MINSUM, Wilcoxon $p=0.002$. Подробный протокол: \texttt{docs/filo2\_inspired\_cap.md}.
  \item[$\square$] Точный объём отчёта: число страниц, таблиц, рисунков, источников и приложений --- уточнить после финальной компиляции PDF (строка реферата).
  \item[$\square$] PNG-визуализации маршрутов для 2--3 инстансов (равномерный, кластерный, выбросы) --- опционально для защиты; JSON с маршрутами доступны в \texttt{data/results/manual\_runs/}.
  \item[$\square$] Полная таблица per-instance результатов --- CSV-файлы доступны в \texttt{data/results/}, можно вынести в отдельное приложение или приложить архивом.
  \item[$\square$] Финальная сверка всех чисел, если проводились дополнительные запуски benchmark после последней правки.
\end{enumerate}

\newpage
\section*{ПРИЛОЖЕНИЕ Б. РЕКОМЕНДУЕМЫЙ СОСТАВ ПОЛНОЙ ТАБЛИЦЫ ЭКСПЕРИМЕНТОВ}
\addcontentsline{toc}{section}{ПРИЛОЖЕНИЕ Б. РЕКОМЕНДУЕМЫЙ СОСТАВ ПОЛНОЙ ТАБЛИЦЫ ЭКСПЕРИМЕНТОВ}

Полную таблицу не следует перегружать в основной части отчета. Ее лучше вынести в приложение или приложить отдельным CSV. Рекомендуемый набор столбцов:

\small
\begin{longtable}{>{\raggedright\arraybackslash}p{0.36\textwidth}>{\raggedright\arraybackslash}p{0.56\textwidth}}
\toprule
\textbf{Столбец} & \textbf{Содержание}\\
\midrule
\endhead
\texttt{instance\_family} & Семейство инстанса: uniform, clustered-center, mixed, outliers и т.д.\\
\texttt{instance} & Имя конкретного файла с тестом.\\
\texttt{node\_count} & Число вершин.\\
\texttt{salesman\_count} & Число агентов.\\
\texttt{solver} & Название алгоритма.\\
\texttt{valid} & Результат проверки корректности маршрутов.\\
\texttt{objective} & Пересчитанное значение MINSUM.\\
\texttt{time\_seconds} & Время выполнения.\\
\texttt{reference\_objective} & Значение внешнего ориентира, если он использовался.\\
\texttt{relative\_gap\_percent} & Относительное отклонение от reference.\\
\texttt{imbalance} & Отношение максимальной длины маршрута к средней длине маршрута.\\
\texttt{seed} & Seed запуска, если алгоритм стохастический.\\
\bottomrule
\end{longtable}
\normalsize

\end{document}
